\documentclass[11pt,letterpaper]{article}

% BEGIN preamble.tex
\usepackage[T1]{fontenc}
\usepackage[utf8]{inputenc}
\usepackage{lmodern}
\usepackage{amsmath,amssymb,amsthm,mathtools,mathrsfs}
\usepackage[margin=1in,headheight=14pt]{geometry}
\usepackage{microtype}
\usepackage{booktabs,longtable,array,enumitem}
\usepackage{xcolor}
\definecolor{accent}{HTML}{184B68}
\definecolor{linkblue}{HTML}{155D8B}
\usepackage{fancyhdr}
\pagestyle{fancy}
\fancyhf{}
\fancyhead[L]{\small\sffamily Tail kernels and boundary products}
\fancyhead[R]{\small\sffamily ProveIt}
\fancyfoot[C]{\thepage}
\renewcommand{\headrulewidth}{0.3pt}
\usepackage{hyperref}
\hypersetup{colorlinks=true,linkcolor=linkblue,citecolor=linkblue,urlcolor=linkblue,
  pdftitle={Bilinear Tail Identities, Polylogarithmic Resonances, and Stieltjes Boundary Products},
  pdfauthor={ProveIt research continuation},bookmarksnumbered=true}
\usepackage{bookmark}
\numberwithin{equation}{section}
\newtheorem{theorem}{Theorem}[section]
\newtheorem{proposition}[theorem]{Proposition}
\newtheorem{lemma}[theorem]{Lemma}
\newtheorem{corollary}[theorem]{Corollary}
\theoremstyle{definition}
\newtheorem{definition}[theorem]{Definition}
\newtheorem{example}[theorem]{Example}
\theoremstyle{remark}
\newtheorem{remark}[theorem]{Remark}
\newtheorem{question}[theorem]{Question}
\DeclareMathOperator{\Li}{Li}
\DeclareMathOperator{\FP}{FP}
\DeclareMathOperator{\CT}{CT}
\DeclareMathOperator{\Res}{Res}
\DeclareMathOperator{\rank}{rank}
\DeclareMathOperator{\diag}{diag}
\DeclareMathOperator{\ord}{ord}
\newcommand{\C}{\mathbb C}
\newcommand{\Q}{\mathbb Q}
\newcommand{\N}{\mathbb N}
\newcommand{\Z}{\mathbb Z}
\newcommand{\Bcal}{\mathcal B}
\newcommand{\fall}[2]{#1^{\underline{#2}}}
\newcommand{\sourcepin}{da2792fc9cc8ef19c6a314be6c65e1437fbdc66b}
\setlength{\parskip}{3pt}
\setlength{\emergencystretch}{2em}
\setcounter{tocdepth}{2}
\allowdisplaybreaks[2]
\setlist{topsep=4pt,itemsep=3pt}

% END preamble.tex

\begin{document}
\hypersetup{pageanchor=false}
\begin{titlepage}
\thispagestyle{empty}
\vspace*{0.6cm}
{\color{accent}\sffamily\large PROVEIT RESEARCH CONTINUATION\par}
\vspace{1.0cm}
{\huge\bfseries Bilinear Tail Identities,\\[5pt]
Polylogarithmic Resonances,\\[5pt]
and Stieltjes Boundary Products\par}
\vspace{0.6cm}
{\Large Complete moment reconstruction at orders ten and twelve\par}
\vspace{0.8cm}
{\large 11 October 2026\par}
\vspace{0.7cm}
\begin{abstract}
We continue the exact-identity programme in the current ProveIt manuscript
and its five incoming research packages. First we classify all rational
certificates formed from the two adjacent products of independently
colored tails. Such a certificate exists precisely when the quotient of
the denominator polynomials is a constant times $h(z)h(z+1)$ with $h$
rational. This gives two closed independently shifted Nielsen-beta
kernels, a rank-two identity, all mixed polygamma derivatives, a Gamma--
Hurwitz rectangular primitive, and a Hermitian polylogarithm-tail law.
Second, a holomorphic combination of two neighboring polylogarithms
evaluates an entire family of Mellin resonances through finite factored
shift operators; its mixed jets have explicit zeta and Stieltjes
coefficients. Third, a completed two-variable Hurwitz product kernel
evaluates every finite-part product having one undifferentiated
Stieltjes factor and one positive argument derivative, including an exact
integration-by-parts endpoint law. Finally, we complete the optimal
moment-prefix conjecture for every unequal even pair whose smaller order
is ten or twelve. The proofs combine exact source reductions, one
five-by-five certificate, an effective analytic threshold, and a fully
replayed finite-field certificate. All convergence conventions and
remaining conjectures are stated explicitly. Historical priority beyond
the inspected literature and repository corpus is not asserted.
\end{abstract}
\vfill
\noindent\textbf{Repository snapshot.}\quad
\texttt{da2792fc9cc8ef19c6a314be6c65e1437fbdc66b}\\[3pt]
\noindent\textbf{Deliverables.}\quad Complete source and PDF, exact verification
programs, recorded certificates, numerical diagnostics, and integration notes.
\end{titlepage}
\pagenumbering{roman}
\hypersetup{pageanchor=true}
\tableofcontents
\clearpage
\pagenumbering{arabic}

% BEGIN sections/01_scope.tex
\section{Research targets, source status, and conventions}
\label{sec:scope}

The relevant source state is the ProveIt commit
\texttt{da2792fc9cc8ef19c6a314be6c65e1437fbdc66b}, dated
11 October 2026, 21:44:07 UTC. We inspected the unified manuscript and
the five archives then present in \texttt{docs/incoming}: the alternating-tail,
finite-difference-determinant, higher-quadratic-jet, higher-spectral-defect,
and quartic-reconstruction reports. Their mathematical status must be read
together. In particular, a question left open in one report can already
have a partial answer in another report written from an earlier snapshot.
The exact paths and selected source passages are recorded in
\texttt{provenance/sources.json}.

The new results concern identities, finite reductions, and exact
reconstruction. The only estimates used below justify interchange of
limits or close a determinant argument; they are not the research target.

The proofs share a concrete ordering principle: enforce cancellation
before summing or specializing parameters. The tail certificates remove
the linear tail terms. The Mellin operators annihilate finite Lerch
corrections. The Hurwitz completions remove the moving coordinate
resonances. The centered moment transform isolates finitely many zeta
coordinates. Keeping the parameters independent until these steps are
complete preserves the mixed derivatives and endpoint information.

\subsection{What advances and what is inherited}

\begin{longtable}{@{}>{\raggedright\arraybackslash}p{0.22\textwidth}>{\raggedright\arraybackslash}p{0.43\textwidth}>{\raggedright\arraybackslash}p{0.28\textwidth}@{}}
\toprule
Source target & Contribution in this article & Scope\\
\midrule\endfirsthead
\toprule Source target & Contribution in this article & Scope\\\midrule\endhead
Two adjacent bilinear tail certificates \cite{tail-alternating}
& Necessary and sufficient rational-factorization criterion; exact
cutoff identity; independent-shift evaluations
& Complete classification within the stated two-product certificate class.\\[5pt]
Independently shifted tail identities
& Two Nielsen-beta kernels, all mixed argument derivatives, and a
Gamma--Hurwitz rectangular primitive
& Two freely varying shifts for adjacent quadratic denominators;
arbitrary independent tail orders are not claimed.\\[5pt]
Polylogarithm derivatives and antiderivatives
& A holomorphic neighboring-order combination, finite Mellin-resonance
operators, and mixed Stieltjes jets
& All complex polylogarithm orders on explicitly stated domains.\\[5pt]
Undifferentiated Stieltjes edge \cite{tail-polygamma}
& Complete joint-pole subtraction and finite ordinary-zeta algorithm
& Every product with exactly one undifferentiated factor; an endpoint
law holds for all derivative orders.\\[5pt]
Optimal first-$q$ moment prefix \cite{tail-alternating,tail-quartic}
& Complete smaller-order-ten and smaller-order-twelve families
& All even $q>10$ and $q>12$, respectively; formal coordinate rank.\\
\bottomrule
\end{longtable}

The bilinear classification and the two completed reconstruction families
answer explicit project questions. The Stieltjes completion answers a
previously excluded boundary case. The Mellin operator and kernel formulas
are proved here and were not located in the inspected project corpus;
their classical ingredients are identified when used. The literature
search is targeted, not an exhaustive priority determination. Products
of zeta tails, Hurwitz product integrals, rational summation, and Mellin
transforms all have substantial prior literature
\cite{Hoffman,EspinosaMoll,ShpotParis,DLMFPolylog}.

\subsection{Notation and analytic conventions}

The Hurwitz zeta function is initially
\[
 \zeta(s,a)=\sum_{n=0}^{\infty}(n+a)^{-s},\qquad
 \Re s>1,\quad\Re a>0,
\]
and is then continued. Its generalized Stieltjes constants satisfy
\begin{equation}\label{eq:stieltjes-convention}
 \zeta(1+u,a)=\frac1u+
 \sum_{m=0}^{\infty}\frac{(-1)^m}{m!}\gamma_m(a)u^m.
\end{equation}
Thus $\gamma_0(a)=-\psi(a)$ and $\gamma_m=\gamma_m(1)$.
A prime on $\zeta(s,a)$ differentiates the order $s$; a superscript
$\gamma_m^{(r)}(a)$ differentiates the argument $a$. We write
\[
 H_n^{(r)}=\sum_{k=1}^n k^{-r},\qquad H_0^{(r)}=0,
 \qquad \psi^{(r)}(a)=(-1)^{r+1}r!\zeta(r+1,a)\quad(r\ge1).
\]
Bernoulli numbers use $B_1=-1/2$ and
$t/(e^t-1)=\sum_{j\ge0}B_jt^j/j!$.
These normalizations agree with the standard sources
\cite{DLMFZeta,DLMFPolygamma}.

All powers and logarithms use the principal branch on the initial slit
or right-half-plane domains. The logarithm of $\Gamma$ on the right
half-plane is the holomorphic branch real on the positive axis.
An ordinary sum is the limit of the displayed partial sums. A finite
part is always specified relative to a coordinate, and a spectral
continuation is kept distinct from an ordinary divergent sum.
Local normal convergence means locally uniform absolute convergence.

The phrase \emph{formal coordinate rank} refers to rational linear
algebra in the explicitly defined source coordinate space. It does not
assert independence of the evaluated periods. The finite-field
computations are exact arithmetic certificates with complete replay;
the analytic arguments remain written mathematical proofs. No
proof-assistant formalization is claimed.

% END sections/01_scope.tex


% BEGIN sections/02_bilinear.tex
\section{Classification of rational bilinear tail certificates}
\label{sec:bilinear}

The latest alternating-tail report asks whether allowing both adjacent
products gives a useful classification for two different denominator
polynomials. We answer that question exactly. The classification is
relative to an explicit certificate class; failure of its criterion
does not exclude longer telescopers or other special-function identities.

\subsection{The certificate class and its factorization criterion}

Let $P,Q\in\C[z]$ be nonzero, let $\alpha,\beta,t\ne0$, and suppose
two sequences satisfy
\begin{equation}\label{bi:recurrences}
 E_{n-1}=E_n+\frac{\alpha^n}{P(n)},\qquad
 F_{n-1}=F_n+\frac{\beta^n}{Q(n)}.
\end{equation}
For rational functions $A,B$, consider
\begin{equation}\label{bi:certificate}
 T_n=t^n\{A(n+1)E_nF_{n+1}+B(n+1)E_{n+1}F_n\}.
\end{equation}
We ask that, after using \eqref{bi:recurrences}, the coefficients of
$E_n$ and $F_n$ in $T_n-T_{n-1}$ both vanish identically. The remaining
term is then a rational weight times $E_nF_n$. This is an algebraic
condition on the certificate, independent of any conjectured numerical
relation among the tails.

\begin{theorem}[Complete two-adjacent-product classification]
\label{bi:classification}
A nonzero rational certificate \eqref{bi:certificate} has the stated
cancellation property if and only if
\begin{equation}\label{bi:criterion}
 t^2\alpha\beta=1,\qquad
 \frac{P(z)}{Q(z)}=\rho h(z)h(z+1)
\end{equation}
for a nonzero constant $\rho$ and a nonzero rational function $h$.
In that case all certificates for the chosen $t$ are scalar multiples
of
\begin{equation}\label{bi:AB}
 A(z)=Q(z)h(z),\qquad B(z)=\kappa Q(z)h(z+1),\qquad
 \kappa=-t\beta.
\end{equation}
Writing $C=A+B$, their exact difference is
\begin{equation}\label{bi:difference}
 T_n-T_{n-1}
   =t^n\{C(n+1)-t^{-1}C(n)\}E_nF_n.
\end{equation}
Thus, when the rational factorization exists, there are exactly two
outer-color choices $t$, and the certificate space for either choice
is one-dimensional.
\end{theorem}

Here ``nonzero'' means that $(A,B)$ is not the zero pair. At a special
parameter value the resulting weight can vanish identically, or the
evaluated tails can make the certificate zero. Such degeneracies do not
change the formal cancellation classification. An example is the
equal-shift limit of the alternating kernel below, whose nontrivial
diagonal identity is recovered by holomorphic continuation.

\begin{proof}
Direct expansion gives the quadratic coefficient
$A(n+1)+B(n+1)-t^{-1}\{A(n)+B(n)\}$. The linear coefficients vanish
exactly when
\begin{equation}\label{bi:linear}
 \beta\frac{A(z+1)}{Q(z+1)}+t^{-1}\frac{B(z)}{Q(z)}=0,
 \qquad
 \alpha\frac{B(z+1)}{P(z+1)}+t^{-1}\frac{A(z)}{P(z)}=0.
\end{equation}
If one of $A,B$ is zero, both are zero. Otherwise set
$h=A/Q$ and $R=P/Q$. The first equation gives
$B=-t\beta Qh(z+1)$; substitution into the second gives
\begin{equation}\label{bi:twostep}
 \frac{R(z+1)}{R(z)}
       =t^2\alpha\beta\frac{h(z+2)}{h(z)}.
\end{equation}
A nonzero rational function $r$ satisfies $r(z+c)/r(z)\to1$ as
$z\to\infty$ for every fixed $c$. Therefore
$t^2\alpha\beta=1$. Equation \eqref{bi:twostep} now says that
$R(z)/\{h(z)h(z+1)\}$ is a rational function with period one.
A nonconstant rational periodic function would have infinitely many
poles, or, if polynomial, could not have a nonzero period. It is
therefore the constant $\rho$. This proves necessity.

Conversely, \eqref{bi:criterion} and \eqref{bi:AB} verify
\eqref{bi:linear} by substitution, proving the difference identity.
If $h_1,h_2$ are two solutions, their quotient $r$ satisfies
$r(z)r(z+1)=\text{constant}$ and hence $r(z+2)=r(z)$.
It is constant by the same rational-periodicity argument. This proves
the uniqueness assertion.
\end{proof}

\subsection{An effective criterion on the divisor of the rational function}

The condition in \eqref{bi:criterion} does not require an infinite
search for $h$. Reduce $R=P/Q$ first, and group its finite zeros and
poles into equivalence classes modulo integer translation. For an orbit
$\xi+\Z$, put
\[
 r_j=\ord_{z=\xi+j}R(z).
\]
Only finitely many $r_j$ are nonzero; poles have negative multiplicity.

\begin{proposition}[Alternating divisor test]
\label{bi:divisor}
The factorization $R(z)=\rho h(z)h(z+1)$ with rational $h$ exists
if and only if, for every such orbit,
\begin{equation}\label{bi:divisorcondition}
 \sum_{j\in\Z}(-1)^j r_j=0.
\end{equation}
When this holds, the multiplicity of $h$ at $\xi+j$ is
\begin{equation}\label{bi:divisorconstruction}
 h_j=\sum_{k\ge0}(-1)^k r_{j+k}.
\end{equation}
The constant $\rho$ then matches the leading coefficients at infinity.
\end{proposition}

\begin{proof}
The multiplicities must satisfy $r_j=h_j+h_{j+1}$.
The unique solution vanishing above the support of $r$ is
\eqref{bi:divisorconstruction}. It also vanishes below that support
exactly when \eqref{bi:divisorcondition} holds. This constructs a
rational $h$ with the required finite divisor. The quotient
$R/(hh(\cdot+1))$ has neither finite zeros nor poles, and its order
at infinity is also zero, because $\sum r_j=2\sum h_j$.
It is therefore a nonzero constant. Changing the orbit representative
only multiplies \eqref{bi:divisorcondition} by a sign.
\end{proof}

This is a finite algebraic test on the reduced rational function. In
particular, its degree at infinity must be even. Generic quadratic
polynomials with four zeros in distinct integer-translation orbits
fail the test. By contrast,
\begin{equation}\label{bi:adjacent-ratio}
 \frac{(z+a)(z+a+1)}{(z+b)(z+b+1)}
       =\frac{z+a}{z+b}\frac{z+a+1}{z+b+1}
\end{equation}
passes for arbitrary independent $a,b$. This example gives the
two-shift identities in the next section. The divisor argument is an
elementary instance of rational difference-equation reasoning; the
contribution here is its exact use to classify this bilinear tail
certificate and all compatible colors.

\subsection{The finite endpoint and the infinite identity}

Assume now that $\deg P=d\ge2$, $\deg Q=e\ge2$, and neither polynomial
vanishes at a positive integer. Let
\[
 0<|\alpha|,|\beta|\le1,\qquad \alpha\ne1,\quad\beta\ne1,
\]
and define the ordinary tails
\begin{equation}\label{bi:tails}
 E_n=\sum_{k>n}\frac{\alpha^k}{P(k)},\qquad
 F_n=\sum_{k>n}\frac{\beta^k}{Q(k)}.
\end{equation}
Suppose the equivalent conditions of Theorem~\ref{bi:classification}
hold. A zero or pole of $h$ at a positive integer would force, by
\eqref{bi:criterion}, alternating nonzero multiplicities at all later
integers, which is impossible for a rational function. Consequently
$h(1)$ is finite and nonzero. The rational functions $A,B$ are regular
at nonnegative integers after removal of any removable singularities.

\begin{theorem}[Bilinear summation with its exact upper endpoint]
\label{bi:sum}
Put $W(n)=C(n+1)-t^{-1}C(n)$ and
\begin{equation}\label{bi:endpoint}
 \mathcal E=
 \beta h(1)E_0-\frac{t^{-1}}{\rho h(1)}F_0
                  -t^{-1}C(0)E_0F_0.
\end{equation}
For every integer $N\ge0$,
\begin{equation}\label{bi:finite}
 \sum_{n=0}^N t^n W(n)E_nF_n=\mathcal E+T_N.
\end{equation}
Moreover,
\begin{equation}\label{bi:infinite}
 \boxed{\displaystyle
 \sum_{n=0}^{\infty}t^n W(n)E_nF_n=\mathcal E,}
\end{equation}
and the series converges absolutely. For fixed $P,Q,h$, it converges
locally normally in the admissible colors away from $1$, with either
locally chosen branch of $t$.
\end{theorem}

\begin{proof}
Sum \eqref{bi:difference} for $1\le n\le N$. Since
$E_1=E_0-\alpha/P(1)$ and $F_1=F_0-\beta/Q(1)$,
\[
 T_0=C(1)E_0F_0-\beta h(1)E_0
                         -\frac{\alpha\kappa}{\rho h(1)}F_0.
\]
Use $\alpha\kappa=-t^{-1}$ and restore the $n=0$ summand.
This gives \eqref{bi:finite}, including the lower-endpoint term
$C(0)E_0F_0$.

Summation by parts, using bounded partial geometric sums, gives
$E_n=O(|\alpha|^{n+1}n^{-d})$ and
$F_n=O(|\beta|^{n+1}n^{-e})$. The bounds are locally uniform as
stated. The factorization forces $\deg h=(d-e)/2$, so $A,B,C$
are $O(n^{(d+e)/2})$. Since $|t|=|\alpha\beta|^{-1/2}$,
the summand and $T_N$ are respectively bounded by a constant times
\[
 |\alpha\beta|^{n/2}n^{-(d+e)/2},\qquad
 |\alpha\beta|^{N/2}N^{-(d+e)/2}.
\]
The first is summable and the second tends to zero. This proves the
infinite identity and the local normal convergence assertion.
\end{proof}

The finite identity remains true outside these sufficient convergence
hypotheses whenever the recurrences make sense. Any additional ordinary
or regularized limit needs its own endpoint analysis.

\subsection{A Hermitian polylogarithm-tail law}

The two inner colors need not agree. Taking conjugate colors gives an
Hermitian identity with explicit dependence on the color.  Its
weights are positive for the power-polynomial specialization below.

\begin{corollary}[Conjugate-color completion]
\label{bi:hermitian}
Let $P$ be a real polynomial of degree at least two with no positive
integer zero. For $0<\theta<2\pi$, write
\[
 E_n(\theta)=\sum_{k>n}\frac{e^{ik\theta}}{P(k)}.
\]
Then
\begin{equation}\label{bi:hermitian-formula}
 \sum_{n\ge0}\{P(n+1)-P(n)\}|E_n(\theta)|^2
 =\cot\frac\theta2\,\Im E_0(\theta)-\Re E_0(\theta)
                         -P(0)|E_0(\theta)|^2.
\end{equation}
For every integer $d\ge2$, in particular,
\begin{equation}\label{bi:polylog-hermitian}
 \boxed{\begin{aligned}
 &\sum_{n\ge0}\{(n+1)^d-n^d\}
 \left|\Li_d(e^{i\theta})-\sum_{k=1}^n\frac{e^{ik\theta}}{k^d}\right|^2\\
 &\hspace{18mm}=\cot\frac\theta2\,\Im\Li_d(e^{i\theta})
                                  -\Re\Li_d(e^{i\theta}).
 \end{aligned}}
\end{equation}
\end{corollary}

\begin{proof}
In Theorem~\ref{bi:sum}, take $Q=P$, $h=1$, $\rho=1$,
$\alpha=e^{i\theta}$, $\beta=e^{-i\theta}$, and $t=1$.
Then $C=(1-e^{-i\theta})P$ and $F_n=\overline{E_n}$.
Divide by $1-e^{-i\theta}$ and use
\[
 \frac{e^{-i\theta}E-\overline E}{1-e^{-i\theta}}
       =-\Re E+\cot(\theta/2)\Im E.
\]
\end{proof}

At the Gaussian color and order two, this gives the explicit identity
\begin{equation}\label{bi:gaussian-example}
 \boxed{\displaystyle
 \sum_{n\ge0}(2n+1)
 \left|\Li_2(i)-\sum_{k=1}^n\frac{i^k}{k^2}\right|^2
          =G+\frac{\pi^2}{48},}
\end{equation}
where $G=\sum_{k\ge0}(-1)^k/(2k+1)^2$ is Catalan's constant.
Here $\Li_2(i)=-\pi^2/48+iG$, obtained by separating even and odd
terms of the defining series. At $\theta=\pi$ the formula recovers
the ordinary-outer alternating square law of \cite{tail-alternating}.
For $d>2$, its limit as $\theta\to0$ is also an ordinary identity:
\[
 \sum_{n\ge0}\{(n+1)^d-n^d\}
                 \{\zeta(d)-H_n^{(d)}\}^2
            =2\zeta(d-1)-\zeta(d).
\]
Indeed the absolute tail is bounded by $O(n^{1-d})$, giving a
summable majorant $O(n^{1-d})$ after multiplication by the weight.
The right side follows by differentiating the color at zero, justified
by $d>2$. This limit is not asserted at $d=2$, where the ordinary
series diverges. The new classification supplies these formulas as
corollaries; they are not separate claims of global literature priority.

% END sections/02_bilinear.tex


% BEGIN sections/03_shifted_kernels.tex
\section{Independent shifts, polygamma kernels, and Gamma primitives}
\label{sec:shifted}

The adjacent-quadratic factorization \eqref{bi:adjacent-ratio} makes
both permitted outer colors explicit. The ordinary-outer branch has
a kernel of rank at most two. The alternating-outer branch has a
divided-difference kernel. Both admit all mixed argument derivatives.

\subsection{Two complete Nielsen-beta kernels}

For $\Re z>0$, define Nielsen's beta function and its adjacent difference
by
\begin{align}
 \Bcal(z)&=\sum_{j\ge0}\frac{(-1)^j}{j+z}
    =\frac12\left\{\psi\left(\frac{z+1}{2}\right)
                           -\psi\left(\frac z2\right)\right\},
       \label{sh:beta}\\
 U(z)&=2\Bcal(z)-\frac1z
     =\sum_{j\ge0}\frac{(-1)^j}{(j+z)(j+z+1)}
     =\int_0^1 x^{z-1}\frac{1-x}{1+x}\,dx.
       \label{sh:U}
\end{align}
The beta series converges by summation by parts; the series defining
$U$ converges absolutely. Its last representation follows from
$\Bcal(z)+\Bcal(z+1)=1/z$. These are standard descriptions of
Nielsen's function \cite{NielsenBeta}; the identities below concern
products of its independently shifted adjacent differences.

\begin{theorem}[Ordinary and alternating outer kernels]
\label{sh:two-kernels}
For $\Re z,\Re w>0$, the following series are locally normally
convergent:
\begin{align}
 K(z,w)&=\sum_{n\ge0}(2n+z+w)U(n+z)U(n+w),\label{sh:Kdef}\\
 S(z,w)&=\sum_{n\ge0}(-1)^nU(n+z)U(n+w).\label{sh:Sdef}
\end{align}
They have the exact evaluations
\begin{equation}\label{sh:K}
 \boxed{\displaystyle
 K(z,w)=\frac{z}{2w}U(z)+\frac{w}{2z}U(w)
                -\left(zw-\frac{z+w}{2}\right)U(z)U(w),}
\end{equation}
and, for $z\ne w$,
\begin{equation}\label{sh:S}
 \boxed{\displaystyle
 S(z,w)=\frac12U(z)U(w)
      +\frac{z^2U(z)-w^2U(w)}{2zw(z-w)}.}
\end{equation}
The apparent singularity of \eqref{sh:S} on $z=w$ is removable:
\begin{equation}\label{sh:Sdiagonal}
 S(z,z)=\frac12U(z)^2+\frac{U(z)}z+\frac12U'(z).
\end{equation}
\end{theorem}

\begin{proof}
Take
\[
 P(k)=(k+z-1)(k+z),\qquad Q(k)=(k+w-1)(k+w),
 \qquad h(k)=\frac{k+z-1}{k+w-1},\quad\rho=1,
\]
and $\alpha=\beta=-1$ in Theorem~\ref{bi:sum}. Then
$E_n=(-1)^{n+1}U(n+z)$ and $F_n=(-1)^{n+1}U(n+w)$.
For $t=1$, the certificate coefficients are
\[
 A(k)=(k+z-1)(k+w),\qquad B(k)=(k+z)(k+w-1).
\]
Hence
\[
 C(k)=2k^2+2(z+w-1)k+2zw-z-w,\qquad
 C(k+1)-C(k)=2(2k+z+w).
\]
The lower endpoint \eqref{bi:endpoint}, divided by two, gives
\eqref{sh:K}. For $t=-1$, change the sign of $B$; now
$C(k)=z-w$ and $C(k+1)+C(k)=2(z-w)$.
The endpoint gives \eqref{sh:S}. Its diagonal limit is
\eqref{sh:Sdiagonal}.

One can also read the exact truncation error from the proof. For the
ordinary kernel it is
\begin{align}\label{sh:finite-endpoint}
 \sum_{n=0}^N(2n+z+w)U(n+z)U(n+w)
 ={}&K(z,w)\notag\\
 &-\frac12\bigl\{(N+z)(N+w+1)U(N+z)U(N+w+1)\notag\\
 &\qquad +(N+z+1)(N+w)U(N+z+1)U(N+w)\bigr\}.
\end{align}
Finally, \eqref{sh:U} gives $U(n+z)=O(n^{-2})$ locally uniformly
in the right half-plane. The ordinary and alternating summands are
therefore $O(n^{-3})$ and $O(n^{-4})$, respectively. This also
justifies the holomorphic diagonal limit.
\end{proof}

For $z=w$, formula \eqref{sh:K} becomes
\[
 K(z,z)=U(z)-z(z-1)U(z)^2,
\]
the adjacent-quadratic square law in \cite{tail-alternating}.
The off-diagonal formula contains two independent shifts; it is not
obtained by setting two parameters inside a single tail and then
squaring that tail.

\subsection{Explicit harmonic and half-shift evaluations}

Let $L=\log2$. The elementary values
\[
 U(1)=2L-1,\qquad U(2)=\frac32-2L,\qquad
 U(1/2)=\pi-2
\]
give
\begin{align}
 \boxed{\sum_{n\ge0}(2n+3)U(n+1)U(n+2)=2(1-L)^2,}
       \label{sh:integer-example}\\
 \boxed{\sum_{n\ge0}(4n+3)U(n+1/2)U(n+1)=(\pi+2)L-2,}
       \label{sh:half-example}\\
 \boxed{\sum_{n\ge0}(-1)^nU(n+1)^2=2L^2-\frac{\pi^2}{12}.}
       \label{sh:alternating-example}
\end{align}
For the last identity use $U'(1)=1-\pi^2/6$ in
\eqref{sh:Sdiagonal}. These are ordinary convergent series.
They can be written entirely with harmonic numbers: if
\[
 H_n^{\mathrm{alt}}=\sum_{k=1}^n\frac{(-1)^{k-1}}k,
 \qquad u_n=2(-1)^n\{L-H_n^{\mathrm{alt}}\}-\frac1{n+1},
\]
then $u_n=U(n+1)$, so \eqref{sh:integer-example} is
$\sum_{n\ge0}(2n+3)u_nu_{n+1}=2(1-L)^2$, while
\eqref{sh:alternating-example} evaluates $\sum_{n\ge0}(-1)^nu_n^2$.
The cancellation in the displayed expression for $u_n$ is part of
the formula; individual unsummed constant pieces are not to be separated.

\subsection{Rank two and all mixed polygamma derivatives}

Put $f(z)=zU(z)$. Formula \eqref{sh:K} has the simpler equivalent form
\begin{equation}\label{sh:ranktwo}
 K(z,w)=f(z)\Bcal(w)+\Bcal(z)f(w)-f(z)f(w).
\end{equation}
Consequently, for any three right-half-plane points $z_1,z_2,z_3$,
\begin{equation}\label{sh:det3}
 \det\left[
 \sum_{n\ge0}(2n+z_i+z_j)U(n+z_i)U(n+z_j)
                  \right]_{1\le i,j\le3}=0.
\end{equation}
More precisely the two-point determinant is
\begin{equation}\label{sh:det2}
 K(z,z)K(w,w)-K(z,w)^2
       =-\{f(z)\Bcal(w)-\Bcal(z)f(w)\}^2.
\end{equation}
These are identities in the complex variables. They follow by writing
the matrix as the product of its two coordinate columns, the fixed
matrix $\left(\begin{smallmatrix}-1&1\\1&0\end{smallmatrix}\right)$,
and the transpose of the coordinate columns. They concern this
parameter-dependent weighted kernel, not the unweighted Gram matrix.

\begin{corollary}[Every pair of argument derivatives]
\label{sh:mixed}
For all integers $r,s\ge0$ and $\Re z,\Re w>0$,
\begin{align}\label{sh:mixed-formula}
 &\sum_{n\ge0}\bigl\{
 (2n+z+w)U^{(r)}(n+z)U^{(s)}(n+w)
 +rU^{(r-1)}(n+z)U^{(s)}(n+w)\notag\\[-2pt]
 &\hspace{43mm}
 +sU^{(r)}(n+z)U^{(s-1)}(n+w)\bigr\}\notag\\
 &\qquad=f^{(r)}(z)\Bcal^{(s)}(w)
       +\Bcal^{(r)}(z)f^{(s)}(w)-f^{(r)}(z)f^{(s)}(w).
\end{align}
A term with zero prefactor is omitted. The corresponding alternating
identity is obtained by applying $\partial_z^r\partial_w^s$ to
\eqref{sh:S}, with its holomorphic diagonal interpretation.
\end{corollary}

\begin{proof}
Differentiating the integral \eqref{sh:U} gives
\[
 U^{(r)}(n+z)=\int_0^1x^{n+z-1}\log^r x\frac{1-x}{1+x}\,dx
                  =O(n^{-r-2})
\]
locally uniformly. This proves the required normal convergence
after every fixed number of derivatives. Apply the product rule to
\eqref{sh:ranktwo}.
\end{proof}

Every term is explicitly a polygamma expression, because
\begin{equation}\label{sh:polygamma}
 U^{(r)}(z)=2^{-r}\left\{
   \psi^{(r)}\left(\frac{z+1}{2}\right)
  -\psi^{(r)}\left(\frac z2\right)\right\}
       +\frac{(-1)^{r+1}r!}{z^{r+1}}.
\end{equation}
For example, $r=s=1$, $z=w=1$ yields
\begin{equation}\label{sh:derivative-example}
 \sum_{n\ge0}\{(2n+2)U'(n+1)^2+2U(n+1)U'(n+1)\}
        =\frac{\pi^2\log2}{3}-4\log^22.
\end{equation}
The two lower-derivative terms in \eqref{sh:mixed-formula} are
required by the derivative of the weight. Omitting them would give
a different, generally false identity.

\subsection{An explicitly anchored Gamma--Hurwitz primitive}

The rank-two form makes a two-variable antiderivative finite as well.
Define the following holomorphic functions on $\Re z>0$:
\begin{align}
 \mathcal L(z)&=\log\Gamma\left(\frac{z+1}{2}\right)
                         -\log\Gamma\left(\frac z2\right),\label{sh:L}\\
 \mathcal Q(z)&=2z\mathcal L(z)-4\left\{
 \zeta'\left(-1,\frac{z+1}{2}\right)
                   -\zeta'\left(-1,\frac z2\right)\right\},\label{sh:Q}\\
 \mathcal U(z)&=2\mathcal L(z)-\log z.\label{sh:primitiveU}
\end{align}
They satisfy
\begin{equation}\label{sh:primitive-derivatives}
 \mathcal L'(z)=\Bcal(z),\qquad
 \mathcal Q'(z)=zU(z),\qquad
 \mathcal U'(z)=U(z).
\end{equation}
For the middle identity, differentiate the Hurwitz shift law:
\[
 \partial_a\zeta'(-1,a)=-\zeta(0,a)+\zeta'(0,a)
        =a-\frac12+\log\Gamma(a)-\frac12\log(2\pi).
\]
It follows that the derivative of the brace in \eqref{sh:Q} is
$1/4+\mathcal L(z)/2$, giving \eqref{sh:primitive-derivatives}.
Only standard Hurwitz special values are used here \cite{DLMFZeta}.

Fix anchors $z_0,w_0$ in the right half-plane. For any holomorphic
$V$, use the notation
\[
 \Delta_zV(n)=V(n+z)-V(n+z_0),\qquad
 \Delta_wV(n)=V(n+w)-V(n+w_0).
\]

\begin{theorem}[Rectangular Gamma--Hurwitz summation identity]
\label{sh:primitive-theorem}
For $z,z_0,w,w_0$ in the right half-plane,
\begin{align}\label{sh:primitive-identity}
 &\sum_{n\ge0}\{\Delta_z\mathcal Q(n)\Delta_w\mathcal U(n)
                 +\Delta_z\mathcal U(n)\Delta_w\mathcal Q(n)\}\notag\\
 &\quad=\Delta_z\mathcal Q(0)\Delta_w\mathcal L(0)
       +\Delta_z\mathcal L(0)\Delta_w\mathcal Q(0)
       -\Delta_z\mathcal Q(0)\Delta_w\mathcal Q(0).
\end{align}
The series is ordinarily and locally normally convergent, and its
value vanishes on either anchor face. Thus \eqref{sh:primitive-identity}
has no unspecified constants of integration.
If $\mathcal R$ denotes its displayed right side, its exact finite
version is
\begin{align}\label{sh:primitive-finite}
 &\sum_{n=0}^{N}\{\Delta_z\mathcal Q(n)\Delta_w\mathcal U(n)
                 +\Delta_z\mathcal U(n)\Delta_w\mathcal Q(n)\}\notag\\
 &\quad=\mathcal R-\frac12\{\Delta_z\mathcal Q(N)\Delta_w\mathcal Q(N+1)
                       +\Delta_z\mathcal Q(N+1)\Delta_w\mathcal Q(N)\}.
\end{align}
\end{theorem}

\begin{proof}
Integrate \eqref{sh:ranktwo} in $z$ from $z_0$ to $z$ and in $w$
from $w_0$ to $w$. In the defining summand,
\[
 (2n+z+w)U(n+z)U(n+w)
   =f(n+z)U(n+w)+U(n+z)f(n+w).
\]
Equations \eqref{sh:primitive-derivatives} give the left side of
\eqref{sh:primitive-identity}; the rank-two expression gives its
right side. The domain is simply connected, so the chosen paths
do not affect the result. Uniformly on compact parameter and anchor
sets, $\Delta_z\mathcal Q(n)=O(n^{-1})$ and
$\Delta_z\mathcal U(n)=O(n^{-2})$. The products are $O(n^{-3})$,
which justifies the integrations and proves the convergence assertion.
Integrating the two separated products in the exact upper endpoint
\eqref{sh:finite-endpoint} gives \eqref{sh:primitive-finite}.
\end{proof}

At equal anchors and equal endpoints, the formula reduces to the
one-parameter identity
\[
 \sum_{n\ge0}\Delta_z\mathcal Q(n)\Delta_z\mathcal U(n)
       =\Delta_z\mathcal Q(0)\Delta_z\mathcal L(0)
                            -\frac12\{\Delta_z\mathcal Q(0)\}^2.
\]
This connects an ordinary series of Gamma-ratio and Hurwitz-zeta
antiderivatives directly to its lower endpoint. The explicit anchor
also makes every further argument derivative unambiguous.

% END sections/03_shifted_kernels.tex


% BEGIN sections/04_mellin.tex
\section{A Mellin transform with exact cancellation at every integer exponent}
\label{mpr:sec:main}

The next family keeps the polylogarithm order complex.  This is essential:
an identity proved only at positive integer orders would not justify the
order derivatives that produce Stieltjes constants below.  The starting
ingredients are the classical Fermi--Dirac integral for the polylogarithm,
the Lerch integral, and Euler's beta integral
\cite{DLMFPolylog,mpr:DLMFLerch,mpr:DLMFbeta}.
The contribution developed here is a finite factored operator for all
rational kernels in the family, together with joint continuation through
the Mellin and polylogarithm singularities.  The full-line family and its
operator were not located in the inspected manuscript or five newest
incoming reports.  This is a statement about the inspected corpus, not a
claim of worldwide priority.  Related work on integrals of products of
polylogarithms includes Freitas and Sofo
\cite{mpr:Freitas,mpr:Sofo}; their quoted finite-interval problems should
not be confused with the full-line transform proved here.

Let
\[
 \Omega=\mathbb C\setminus(-\infty,0],\qquad
 (\mathsf Ef)(s)=f(s-1),\qquad
 \mathscr Z(s)=s\zeta(s+1).
\]
The function $\mathscr Z$ is entire, with $\mathscr Z(0)=1$.
Products such as $(\mathsf E)_k$ are ordinary polynomials in the shift
operator, using the rising-factorial convention; in particular
$(\mathsf E)_2=\mathsf E(\mathsf E+1)$.
For $m\ge1$ define
\begin{equation}\label{mpr:eq:integral}
 \mathcal I_m(t,s;\lambda)
 =\int_0^\infty
       \frac{x^{t-1}\operatorname{Li}_s(-\lambda x)}{(1+x)^m}\,dx,
 \qquad -1<\Re t<m,\quad s\in\mathbb C,\quad\lambda\in\Omega.
\end{equation}
Here $x^{t-1}=\exp((t-1)\log x)$ uses the real logarithm, while the
polylogarithm uses its principal argument branch.  Thus its argument never
crosses $[1,\infty)$ for the displayed values of $\lambda$.
The stated $t$ strip is a common convergence domain for every order.
Special orders, notably negative integers, can have larger convergence
strips; no maximal-strip assertion for each individual order is intended.

\subsection{Analyticity of the ordinary integrals}

\begin{lemma}\label{mpr:lem:holomorphic}
The integral \eqref{mpr:eq:integral} is absolutely convergent and jointly
holomorphic in its three variables on the displayed domain.  Every fixed
number of derivatives with respect to $s$, $t$, and $\lambda$ may be taken
under the integral sign.  A $t$ derivative inserts a factor $\log x$, and
$\lambda\partial_\lambda$ replaces the order $s$ by $s-1$.
\end{lemma}

\begin{proof}
We record a compact-set bound that also covers negative and noninteger
orders.  If $s$ and $\lambda$ range over compact subsets of $\mathbb C$
and $\Omega$, respectively, then
\[
 \operatorname{Li}_s(-\lambda x)=O(x)\quad(x\downarrow0),\qquad
 \operatorname{Li}_s(-\lambda x)=O((1+\log x)^B)\quad(x\ge1)
\]
for a sufficiently large fixed $B$.  The first bound follows from the
normally convergent power series near zero.  For the second, choose an
integer $N$ with $\Re(s+N)>0$ throughout a slightly larger compact set and
use
\[
 \operatorname{Li}_{s+N}(-\lambda x)
 =-\frac1{\Gamma(s+N)}\int_0^\infty
        u^{s+N-1}\frac{\lambda x}{e^u+\lambda x}\,du.
\]
Apply $(x\partial_x)^N$.  For a complex variable on rays uniformly
separated from the negative real axis, every fixed logarithmic derivative
of $w/(1+w)$ is bounded by $C\min(|w|,1)$.  Splitting the $u$ integral
at $\log x+C_0$ therefore gives the asserted logarithmic bound.
The identity $(x\partial_x)^N\operatorname{Li}_{s+N}
=\operatorname{Li}_s$ holds first by the power series and then by analytic
continuation.  Cauchy's formula on the slightly larger compact parameter
sets supplies the same sort of bounds for all fixed parameter derivatives.

The resulting majorants at the two endpoints are
$x^{\Re t}$ and $x^{\Re t-m-1}(1+\log x)^B$.  They are locally uniformly
integrable throughout the common strip used in
\eqref{mpr:eq:integral}.  Inserting any fixed power of $\log x$ preserves
integrability.  The standard dominated holomorphy theorem completes the
proof.
\end{proof}

\subsection{The paired Lerch expression}

For $\Re a>0$, first on $|\lambda|<1$ within $\Omega$, put
\begin{equation}\label{mpr:eq:Hdefinition}
 \mathcal H_a(s,\lambda)
 =\operatorname{Li}_s(\lambda)
        -\lambda^a\Phi(\lambda,s,a),
 \qquad \lambda^a=\exp(a\operatorname{Log}\lambda).
\end{equation}
Its two terms must be continued together.  This convention has a concrete
ordinary-integral explanation.

\begin{lemma}\label{mpr:lem:H}
The function \eqref{mpr:eq:Hdefinition} extends jointly holomorphically
to $s\in\mathbb C$, $\Re a>0$, and $\lambda\in\Omega$.  For $\Re s>0$,
\begin{equation}\label{mpr:eq:Hkernel}
 \mathcal H_a(s,\lambda)
 =\frac1{\Gamma(s)}\int_0^\infty u^{s-1}
       \frac{w-w^a}{1-w}\,du,
 \qquad w=\lambda e^{-u},\quad
 w^a=\exp(a(\operatorname{Log}\lambda-u)).
\end{equation}
The apparent singularity at $w=1$ has the removable value $a-1$.
Consequently
\begin{align}
 \mathcal H_a(s,1)&=\zeta(s)-\zeta(s,a),
       \label{mpr:eq:Hone}\\
 \mathcal H_M(s,\lambda)&=\sum_{n=1}^{M-1}\frac{\lambda^n}{n^s}
       \quad(M=1,2,\ldots).
       \label{mpr:eq:Hinteger}
\end{align}
The pole of each zeta function in \eqref{mpr:eq:Hone} cancels.
\end{lemma}

\begin{proof}
For $|\lambda|<1$ and $\Re s>0$, the defining geometric series and the
gamma integral give \eqref{mpr:eq:Hkernel}.  The quotient is analytic at
$w=1$ because its numerator and denominator vanish to first order and
their ratio tends to $a-1$.  It and all fixed $u$ derivatives decay
exponentially as $u\to\infty$, locally uniformly in $a$ and $\lambda$.
They are bounded near $u=0$, including when $\lambda=1$.

To obtain the continuation in order without appealing to a divergent
Mellin integral, let $K_a(u,\lambda)=(w-w^a)/(1-w)$.  For any integer
$N\ge0$, integration by parts first on $\Re s>0$ gives
\[
 \mathcal H_a(s,\lambda)
 =\frac1{\Gamma(s+N)}\int_0^\infty
        u^{s+N-1}(-\partial_u)^N K_a(u,\lambda)\,du.
\]
The right side is holomorphic for $\Re s>-N$.  All boundary terms vanish
in the initial domain; the formulas agree on overlaps by continuation.
Since $N$ is arbitrary, they establish entire dependence on $s$ and the
claimed joint holomorphy.  Substitution of $\lambda=1$ proves
\eqref{mpr:eq:Hone} first where the separate zeta integrals converge and
then everywhere.  For a positive integer $M$, the kernel is the finite
geometric sum $w+\cdots+w^{M-1}$, proving
\eqref{mpr:eq:Hinteger}.
\end{proof}

\subsection{The full transform and the cancellation polynomial}

Define the explicit polynomial
\begin{equation}\label{mpr:eq:P}
 P_{m,t}(X)=(1-t-X)_{m-1}
           =\sum_{j=0}^{m-1}c_{m,j}(t)X^j.
\end{equation}

\begin{theorem}[Full Mellin transform]\label{mpr:thm:master}
For every parameter in \eqref{mpr:eq:integral},
\begin{equation}\label{mpr:eq:master}
 \boxed{\displaystyle
 \mathcal I_m(t,s;\lambda)
 =\frac{\pi}{(m-1)!\sin\pi t}
       \sum_{j=0}^{m-1}c_{m,j}(t)
                  \mathcal H_{m-t}(s-j,\lambda).}
\end{equation}
At $t=0,1,\ldots,m-1$ the numerator vanishes identically as a function
of $s$ and $\lambda$; the apparent poles are removable.  At $\lambda=1$
this becomes the entire-order Hurwitz formula
\begin{equation}\label{mpr:eq:hurwitz}
 \mathcal I_m(t,s;1)
 =\frac{\pi}{(m-1)!\sin\pi t}
       \sum_{j=0}^{m-1}c_{m,j}(t)
       \bigl(\zeta(s-j)-\zeta(s-j,m-t)\bigr).
\end{equation}
\end{theorem}

\begin{proof}
First take $m=1$, $0<\lambda<1$, $\Re s>0$, and $-1<\Re t<0$.
Insert the Fermi--Dirac integral in Lemma~\ref{mpr:lem:holomorphic}.
Absolute convergence permits interchanging the two integrations.  The
elementary partial fraction integral
\[
 \int_0^\infty\frac{x^t}{(1+x)(x+c)}\,dx
  =-\frac{\pi}{\sin\pi t}\frac{1-c^t}{c-1},\qquad c>1,
\]
follows from the beta integral.  With $c=e^u/\lambda$, it gives
\[
 \mathcal I_1(t,s;\lambda)
 =\frac\pi{\sin\pi t}
       \bigl(\operatorname{Li}_s(\lambda)
           -\lambda^{1-t}\Phi(\lambda,s,1-t)\bigr).
\]
The preceding two lemmas continue this identity in $s$ and $\lambda$,
and across the strip $-1<\Re t<1$ after removing $t=0$.

For $m\ge2$, integrate the derivative of
$x^t(1+x)^{1-m}\operatorname{Li}_s(-\lambda x)$ over the positive
axis.  On the common strip $-1<\Re t<m-1$ its boundary terms vanish.
Using $\mathcal I_m(t+1,s;\lambda)
=\mathcal I_{m-1}(t,s;\lambda)-\mathcal I_m(t,s;\lambda)$ yields
\begin{equation}\label{mpr:eq:recurrence}
 \mathcal I_m(t,s;\lambda)
 =\frac{m-1-t-\mathsf E}{m-1}\,
                     \mathcal I_{m-1}(t,s;\lambda).
\end{equation}
Iteration gives $P_{m,t}(\mathsf E)\mathcal I_1/(m-1)!$.

Initially on $|\lambda|<1$, the elementary Lerch shift relation gives
\[
 \mathcal H_{1-t}(s,\lambda)-\mathcal H_{m-t}(s,\lambda)
   =-\sum_{n=0}^{m-2}
       \frac{\lambda^{n+1-t}}{(n+1-t)^s}.
\]
For every summand, $\mathsf E$ acts as multiplication by $n+1-t$;
$P_{m,t}(n+1-t)=(-n)_{m-1}=0$.  Applying the polynomial therefore
annihilates the entire finite shift.  This proves
\eqref{mpr:eq:master} on a nonempty open strip, with the important
positive-real-part argument $m-t$.

At $t=k\in\{0,\ldots,m-1\}$, equation
\eqref{mpr:eq:Hinteger} says that its numerator equals
\[
 \sum_{n=1}^{m-k-1}\frac{\lambda^n}{n^s}
                    (1-k-n)_{m-1}=0.
\]
Every product in the last sum has a zero factor.  Thus all apparent
poles inside $-1<\Re t<m$ are removable.  The two sides are
holomorphic on that connected strip and agree on an open subset, proving
the full domain.  Equation \eqref{mpr:eq:Hone} gives the specialization
\eqref{mpr:eq:hurwitz}.
\end{proof}

For example, the base specialization is
\begin{equation}\label{mpr:eq:base}
 \int_0^\infty\frac{x^{t-1}\operatorname{Li}_s(-x)}{1+x}\,dx
  =\frac\pi{\sin\pi t}\bigl(\zeta(s)-\zeta(s,1-t)\bigr),
 \qquad -1<\Re t<1.
\end{equation}
The difference on the right is the generalized harmonic function
$H_{-t}^{(s)}$.  Its singularities at special orders have cancelled
before order differentiation is performed.

\section{The branch cancellation and a factored formula at resonance}
\label{mpr:sec:resonance}

For $|\lambda|<1$ in $\Omega$, define
\begin{equation}\label{mpr:eq:A}
 \mathcal A_s(\lambda)
    =\operatorname{Log}\lambda\,\operatorname{Li}_s(\lambda)
                         -s\operatorname{Li}_{s+1}(\lambda).
\end{equation}

\begin{theorem}[Continuation through the positive polylogarithm cut]
\label{mpr:thm:A}
The paired expression \eqref{mpr:eq:A} extends jointly holomorphically
to $s\in\mathbb C$, $\lambda\in\Omega$.  It satisfies
\begin{align}
 \mathcal A_s(\lambda)
    &=\int_0^\infty
        \frac{\operatorname{Li}_s(-\lambda x)}{x(1+x)}\,dx,
          \label{mpr:eq:Aintegral}\\
 \mathcal A_s(1)&=-\mathscr Z(s),
          \label{mpr:eq:Aone}\\
 \lambda\partial_\lambda\mathcal A_s(\lambda)
          &=\mathcal A_{s-1}(\lambda).
          \label{mpr:eq:Aderivative}
\end{align}
On the real interval $\lambda>1$, either consistent boundary value of
the two principal polylogarithms in \eqref{mpr:eq:A} gives the same
answer.  At $\lambda=1$ the full expression is removable even at
integer orders where its separate gamma and zeta expansions have poles.

If $u=\operatorname{Log}\lambda$ and $|u|<2\pi$, then
\begin{equation}\label{mpr:eq:localexpansion}
 \boxed{\displaystyle
 \mathcal A_s(e^u)
      =-\sum_{n=0}^\infty\mathscr Z(s-n)\frac{u^n}{n!}.}
\end{equation}
The series converges normally for $s$ in compact subsets of $\mathbb C$
and $|u|<2\pi$.  On the full $u$ disk it describes the continuation of
the germ at zero.  For the original principal composition one retains
$|\Im u|<\pi$, as is automatic when $u=\operatorname{Log}\lambda$.
\end{theorem}

\begin{proof}
Take the removable limit $t\to0$ in the $m=1$ transform, initially for
$|\lambda|<1$.  The Lerch derivative
$\partial_a\Phi(\lambda,s,a)=-s\Phi(\lambda,s+1,a)$ gives
\[
 \left.\partial_t
       \bigl(\operatorname{Li}_s(\lambda)
             -\lambda^{1-t}\Phi(\lambda,s,1-t)\bigr)\right|_{t=0}
 =\operatorname{Log}\lambda\,\operatorname{Li}_s(\lambda)
                         -s\operatorname{Li}_{s+1}(\lambda).
\]
This proves \eqref{mpr:eq:Aintegral}; the integral supplies the claimed
joint holomorphic continuation.  In particular, it proves that the two
boundary values across the positive cut agree, without assigning separate
divergent values.  For $\Re s>0$, taking $\lambda\to1$ from below in
\eqref{mpr:eq:A} gives $-s\zeta(s+1)$.  Both sides are entire in $s$,
proving \eqref{mpr:eq:Aone} at every order.  Differentiation of
\eqref{mpr:eq:Aintegral} proves \eqref{mpr:eq:Aderivative}.

There is also an explicit check of the branch cancellation.  For
noninteger $s$, the classical local polylogarithm expansion
\cite[Eq.~25.12.12]{DLMFPolylog} is
\[
 \operatorname{Li}_s(e^u)
  =\Gamma(1-s)(-u)^{s-1}
      +\sum_{n\ge0}\zeta(s-n)\frac{u^n}{n!},\qquad |u|<2\pi,
\]
on an initial compatible branch.  The two fractional-power terms in
$u\operatorname{Li}_s(e^u)-s\operatorname{Li}_{s+1}(e^u)$ cancel
because $\Gamma(1-s)=-s\Gamma(-s)$.  Combining the remaining series
gives precisely \eqref{mpr:eq:localexpansion}.

Every coefficient $\mathscr Z(s-n)$ is entire in $s$.  The zeta
functional equation and the gamma-ratio estimate imply, uniformly on a
compact set of orders,
\[
 \frac{|\mathscr Z(s-n)|}{n!}
        \le C_K n^{B_K}(2\pi)^{-n}
\]
for all sufficiently large $n$.  Consequently the series is normally
convergent on the claimed product domain.  Continuation in $s$ now
includes all integer orders and proves the claimed cancellation there.
The restriction on the original principal logarithm is stated explicitly
to avoid confusing a germ in $u$ with a globally single-valued expression
under the exponential map.
\end{proof}

\begin{theorem}[Every integer Mellin resonance]\label{mpr:thm:integer}
For $m\ge1$, $0\le k\le m-1$, $s\in\mathbb C$, and $\lambda\in\Omega$,
\begin{equation}\label{mpr:eq:integer}
 \boxed{\displaystyle
 \int_0^\infty\frac{x^{k-1}\operatorname{Li}_s(-\lambda x)}{(1+x)^m}\,dx
 =\frac{(\mathsf E)_k(1-\mathsf E)_{m-k-1}}{(m-1)!}
                                  \mathcal A_s(\lambda).}
\end{equation}
In particular,
\begin{equation}\label{mpr:eq:integerzeta}
 \int_0^\infty\frac{x^{k-1}\operatorname{Li}_s(-x)}{(1+x)^m}\,dx
 =-\frac{(\mathsf E)_k(1-\mathsf E)_{m-k-1}}{(m-1)!}
                                  \mathscr Z(s).
\end{equation}
The right side is a finite expression in entire functions.  Thus every
fixed number of order derivatives is valid, including at the values
where a displayed individual zeta factor has a pole.
\end{theorem}

\begin{proof}
At $t=0$, recurrence \eqref{mpr:eq:recurrence} and
\eqref{mpr:eq:Aintegral} give
$\mathcal I_m(0,s;\lambda)
=(1-\mathsf E)_{m-1}\mathcal A_s(\lambda)/(m-1)!$.
For $k\ge1$, multiplication of the kernel by
$x=(1+x)-1$ gives
\[
 \mathcal I_m(k,s;\lambda)
  =\mathcal I_{m-1}(k-1,s;\lambda)
                -\mathcal I_m(k-1,s;\lambda).
\]
Induct on $k$.  The claimed polynomials satisfy exactly this recurrence:
\begin{align*}
 &\frac{(X)_{k-1}(1-X)_{m-k-1}}{(m-2)!}
  -\frac{(X)_{k-1}(1-X)_{m-k}}{(m-1)!}\\
 &\hspace{15mm}
 =\frac{(X)_{k-1}(1-X)_{m-k-1}}{(m-1)!}
                   \bigl((m-1)-(m-k-X)\bigr)\\
 &\hspace{15mm}
 =\frac{(X)_k(1-X)_{m-k-1}}{(m-1)!}.
\end{align*}
This proves \eqref{mpr:eq:integer}.  Equation
\eqref{mpr:eq:Aone} proves \eqref{mpr:eq:integerzeta}.
\end{proof}

The first few kernels illustrate the factored operator:
\[
\begin{array}{c|c}
 (m,k)&(m-1)!^{-1}(X)_k(1-X)_{m-k-1}\\ \hline
 (1,0)&1\\
 (2,0)&1-X\\
 (2,1)&X\\
 (3,1)&(X-X^2)/2\\
 (3,2)&(X+X^2)/2\\
 (4,2)&(X-X^3)/6
\end{array}
\]
For instance, uniformly as an entire identity in $s$,
\begin{align}
 \int_0^\infty\frac{\operatorname{Li}_s(-x)}{(1+x)^2}\,dx
       &=-(s-1)\zeta(s),\label{mpr:eq:squared}\\
 \int_0^\infty\frac{x\operatorname{Li}_s(-x)}{(1+x)^4}\,dx
       &=\frac{(s-3)\zeta(s-2)-(s-1)\zeta(s)}6.
       \label{mpr:eq:quartic}
\end{align}
In \eqref{mpr:eq:squared}, the value at $s=1$ is $-1$, not zero;
in \eqref{mpr:eq:quartic}, the value of $(s-3)\zeta(s-2)$ at $s=3$
is $1$.  These are ordinary removable values of entire products.

\subsection{A bilinear identity isolating a single Stieltjes constant}

The same transform has a particularly short bilinear consequence.
One polylogarithm has a fixed negative integer order and is therefore
a rational function; the other order remains an arbitrary complex
variable.  No relation is differentiated outside its proven domain.

\begin{corollary}\label{mpr:cor:bilinear}
For every integer $d\ge1$, every $s\in\mathbb C$, and every
$\lambda\in\Omega$,
\begin{equation}\label{mpr:eq:bilinear}
 \boxed{\displaystyle
 \int_0^\infty\operatorname{Li}_s(-\lambda x)
                    \operatorname{Li}_{-d}(-x)\,\frac{dx}{x}
   =(-1)^d\mathcal A_{s-d}(\lambda).}
\end{equation}
The integral is ordinary and absolutely convergent, and is entire in
$s$.  At $\lambda=1$ it equals
$(-1)^{d+1}(s-d)\zeta(s-d+1)$, with the removable value understood.
Consequently, for every $r\ge1$,
\begin{equation}\label{mpr:eq:single-stieltjes}
 \boxed{\displaystyle
 \int_0^\infty\operatorname{Li}_{-d}(-x)
       \left.\partial_s^r\operatorname{Li}_s(-x)\right|_{s=d}
                 \frac{dx}{x}
        =(-1)^{d+r}r\gamma_{r-1}.}
\end{equation}
At derivative order zero the integral is $(-1)^{d+1}$.
\end{corollary}

\begin{proof}
Put $D=x\partial_x$ and $f(x)=1/(1+x)$.  Since
$f=1+\operatorname{Li}_0(-x)$, one has
$D^df=\operatorname{Li}_{-d}(-x)$ for $d\ge1$.
Every $D^jf$ is $O(x^{-1})$ at infinity.  At zero it is bounded
when $j=0$ and $O(x)$ when $j\ge1$.
Lemma~\ref{mpr:lem:holomorphic} therefore justifies $d$ successive
integrations by parts with respect to $dx/x$; every boundary term
vanishes.  They give
\[
 \int_0^\infty\operatorname{Li}_s(-\lambda x)D^df(x)\,\frac{dx}{x}
  =(-1)^d\int_0^\infty
       \frac{\operatorname{Li}_{s-d}(-\lambda x)}{x(1+x)}\,dx.
\]
Equation~\eqref{mpr:eq:Aintegral} proves \eqref{mpr:eq:bilinear}.
The same endpoint bounds give normal convergence in order, so the
order derivatives are valid.  At $\lambda=1$, use
$\mathcal A_v(1)=-\mathscr Z(v)$ and
$\mathscr Z^{(r)}(0)=r(-1)^{r-1}\gamma_{r-1}$ to obtain
\eqref{mpr:eq:single-stieltjes}.
\end{proof}

For example the choices $d=1$ and $d=2$ use only the kernels
\[
 \frac{\operatorname{Li}_{-1}(-x)}x=-\frac1{(1+x)^2},\qquad
 \frac{\operatorname{Li}_{-2}(-x)}x=\frac{x-1}{(1+x)^3}.
\]
Thus arbitrary Stieltjes indices occur in ordinary integrals with
fixed rational kernels.  The identity is a consequence of the
proved entire-order transform; no arithmetic independence assertion
is involved.

\section{All mixed logarithmic and Stieltjes identities in the family}
\label{mpr:sec:jets}

For $a\notin\{0,-1,-2,\ldots\}$ with the usual compatible branch,
write
\[
 \zeta(1+v,a)=\frac1v+
       \sum_{r=0}^\infty\frac{(-1)^r}{r!}\gamma_r(a)v^r,
 \qquad \gamma_r=\gamma_r(1).
\]
Only $\Re a>0$ is needed in the formulas below.

\begin{corollary}[Generalized Stieltjes differences]
\label{mpr:cor:stieltjes}
For $r\ge0$ and $0<\Re a<2$,
\begin{equation}\label{mpr:eq:generalizedstieltjes}
 \int_0^\infty\frac{x^{-a}}{1+x}
   \left.\partial_s^r\operatorname{Li}_s(-x)\right|_{s=1}\,dx
  =\frac{(-1)^r\pi}{\sin\pi a}
                     \bigl(\gamma_r-\gamma_r(a)\bigr).
\end{equation}
The value at $a=1$ is removable.  More explicitly it is
\begin{equation}\label{mpr:eq:aone}
 -\zeta^{(r)}(2)-r\zeta^{(r-1)}(2),
\end{equation}
where the second term is omitted for $r=0$.
\end{corollary}

\begin{proof}
Set $t=1-a$ in \eqref{mpr:eq:base} and differentiate the entire
difference $\zeta(s)-\zeta(s,a)$ $r$ times at $s=1$.
The common pole cancels before the coefficient is taken.  Lemma
\ref{mpr:lem:holomorphic} justifies the differentiated integral.
At $a=1$, use \eqref{mpr:eq:Aone} and differentiate $s\zeta(s+1)$ at
$s=1$, obtaining \eqref{mpr:eq:aone}.
\end{proof}

For a particularly explicit example, the half-argument distribution
formula $\zeta(s,1/2)=(2^s-1)\zeta(s)$ gives
\begin{equation}\label{mpr:eq:halfexample}
 \int_0^\infty\frac{x^{-1/2}}{1+x}
   \left.\partial_s\operatorname{Li}_s(-x)\right|_{s=1}\,dx
 =-\pi\bigl(2\gamma\log2+\log^2 2\bigr).
\end{equation}

\begin{corollary}[A single Stieltjes constant at every order]
\label{mpr:cor:quarticstieltjes}
For every $r\ge1$,
\begin{equation}\label{mpr:eq:quarticstieltjes}
 \boxed{\displaystyle
 \int_0^\infty\frac{x}{(1+x)^4}
    \left.\partial_s^r\operatorname{Li}_s(-x)\right|_{s=3}\,dx
 =\frac{r(-1)^{r-1}\gamma_{r-1}
       -2\zeta^{(r)}(3)-r\zeta^{(r-1)}(3)}6.}
\end{equation}
At order zero the same integral equals $(1-2\zeta(3))/6$.
\end{corollary}

\begin{proof}
The entire normalization satisfies
\[
 \mathscr Z(v)=1+\sum_{r\ge1}
                 \frac{(-1)^{r-1}\gamma_{r-1}}{(r-1)!}v^r,
 \qquad
 \mathscr Z^{(r)}(0)=r(-1)^{r-1}\gamma_{r-1}.
\]
Differentiate \eqref{mpr:eq:quartic} at $s=3$.  Its other entire factor
has derivative
$\mathscr Z^{(r)}(2)=2\zeta^{(r)}(3)+r\zeta^{(r-1)}(3)$.
This proves the formula, including the stated order-zero value.
\end{proof}

Thus the first two derivatives give the concrete ordinary integrals
\begin{align*}
 \int_0^\infty\frac{x}{(1+x)^4}
    \left.\partial_s\operatorname{Li}_s(-x)\right|_{s=3}\,dx
   &=\frac{\gamma-\zeta(3)-2\zeta'(3)}6,\\
 \int_0^\infty\frac{x}{(1+x)^4}
    \left.\partial_s^2\operatorname{Li}_s(-x)\right|_{s=3}\,dx
   &=-\frac{\gamma_1+\zeta'(3)+\zeta''(3)}3.
\end{align*}

\subsection{Every power of the logarithmic integration variable}

Write
\[
 \frac{\pi h}{\sin\pi h}=\sum_{a\ge0}b_{2a}h^{2a},\qquad
 b_0=1,\qquad b_{2a}=2(1-2^{1-2a})\zeta(2a)\quad(a\ge1).
\]
For example $b_2=\zeta(2)$ and $b_4=7\zeta(4)/4$.

\begin{theorem}[All logarithmic moments]\label{mpr:thm:logmoments}
For every integer $\ell\ge0$ and every complex $s$,
\begin{equation}\label{mpr:eq:logmoments}
 \boxed{\displaystyle
 \int_0^\infty
       \frac{\log^\ell x\operatorname{Li}_s(-x)}{x(1+x)}\,dx
 =-\ell!\sum_{a=0}^{\lfloor\ell/2\rfloor}
        b_{2a}\frac{(s)_{\ell+1-2a}}{(\ell+1-2a)!}
                    \zeta(s+\ell+1-2a).}
\end{equation}
Every product $(s)_n\zeta(s+n)$ occurring here is entire, with its
removable value understood.  Every fixed number of order derivatives
can be taken on both sides.
\end{theorem}

\begin{proof}
The Hurwitz Taylor series at $a=1$, obtained by repeated use of
$\partial_a\zeta(s,a)=-s\zeta(s+1,a)$, gives
\[
 \zeta(s)-\zeta(s,1-h)
    =-\sum_{n\ge1}\frac{(s)_n}{n!}\zeta(s+n)h^n,
       \qquad |h|<1.
\]
The putative pole at $s=1-n$ cancels against the corresponding linear
factor of $(s)_n$.  The identity holds normally for $s$ in compact sets
by the holomorphy of the Hurwitz difference and Cauchy's formula.
Multiply this series by $\pi/\sin\pi h$, take the coefficient of
$h^\ell$, and apply Lemma~\ref{mpr:lem:holomorphic} to
\eqref{mpr:eq:base}.  This proves the formula and its order derivatives.
\end{proof}

The first two nonconstant logarithmic moments are
\begin{align*}
 \int_0^\infty\frac{\log x\operatorname{Li}_s(-x)}{x(1+x)}\,dx
   &=-\frac{s(s+1)}2\zeta(s+2),\\
 \int_0^\infty\frac{\log^2x\operatorname{Li}_s(-x)}{x(1+x)}\,dx
   &=-\frac{s(s+1)(s+2)}3\zeta(s+3)
                           -2\zeta(2)s\zeta(s+1).
\end{align*}
In particular,
\[
 \int_0^\infty\frac{\log x}{x(1+x)}
      \left.\partial_s\operatorname{Li}_s(-x)\right|_{s=-1}\,dx
       =\frac{\gamma-1}{2}.
\]
This follows by expanding
$-s(s+1)\zeta(s+2)/2$ at $s=-1$; its value before differentiating is
$1/2$.

For a general kernel and general scale, all Mellin derivatives at a
resonant exponent have an equally finite coefficient prescription.
Let
\[
 N_m(t,s,\lambda)
     =\sum_{j=0}^{m-1}c_{m,j}(t)\mathcal H_{m-t}(s-j,\lambda).
\]
Since $N_m(k,s,\lambda)=0$, Theorem~\ref{mpr:thm:master} gives
\begin{equation}\label{mpr:eq:mixedcompiler}
 \left.\partial_t^\ell\partial_s^r
                      \mathcal I_m(t,s;\lambda)\right|_{t=k}
 =\frac{(-1)^k\ell!}{(m-1)!}
   [h^{\ell+1}]\frac{\pi h}{\sin\pi h}\,
                       \partial_s^r N_m(k+h,s,\lambda).
\end{equation}
Only finitely many derivatives and coefficients are required for each
fixed $\ell,r,m$.  This formula also retains mixed $s$ and $\lambda$
derivatives, since $\lambda\partial_\lambda$ simply shifts the order.
It specifies the transverse coordinate $t-k$ and never differentiates a
formula that was known only on the resonant set.

\section{Exact scaled examples and further directions}
\label{mpr:sec:examples}

Put $L=\log2$.  Classical values at $1/2$ yield
\begin{align}
 \mathcal A_1(\lambda)
       &=\operatorname{Li}_2(1-\lambda)-\zeta(2),
            \label{mpr:eq:A1}\\
 \mathcal A_2(1/2)
       &=-\frac74\zeta(3)+\frac{\pi^2L}{12}+
                                               \frac{L^3}{6},
            \label{mpr:eq:Ahalf}\\
 \mathcal A_2(2)
       &=-\frac74\zeta(3)-\frac{\pi^2L}{4}.
            \label{mpr:eq:Atwo}
\end{align}
Equation \eqref{mpr:eq:A1} is Euler's dilogarithm reflection identity in
the present normalization.  Equation \eqref{mpr:eq:Ahalf} follows from
the standard formulas for $\operatorname{Li}_2(1/2)$ and
$\operatorname{Li}_3(1/2)$.
For \eqref{mpr:eq:Atwo}, the compatible boundary values are
\[
 \operatorname{Li}_2(2)=\frac{\pi^2}{4}-i\pi L,\qquad
 \operatorname{Li}_3(2)=\frac78\zeta(3)+\frac{\pi^2L}{4}
                                           -\frac{i\pi L^2}{2}.
\]
They follow from classical inversion and the half-argument evaluations;
their imaginary parts cancel in $L\operatorname{Li}_2(2)
-2\operatorname{Li}_3(2)$.  The opposite compatible boundary convention
gives their conjugates and the same real answer.

Theorem~\ref{mpr:thm:integer} consequently gives, without any integer
relation search,
\begin{align}
 \int_0^\infty\frac{\operatorname{Li}_2(-x/2)}{x(1+x)}\,dx
       &=-\frac74\zeta(3)+\frac{\pi^2L}{12}+\frac{L^3}{6},
             \label{mpr:eq:ex1}\\
 \int_0^\infty\frac{\operatorname{Li}_2(-2x)}{x(1+x)}\,dx
       &=-\frac74\zeta(3)-\frac{\pi^2L}{4},
             \label{mpr:eq:ex2}\\
 \int_0^\infty\frac{\operatorname{Li}_3(-2x)}{(1+x)^3}\,dx
       &=\frac{\pi^2(1-L)}8-\frac78\zeta(3).
             \label{mpr:eq:ex3}
\end{align}
All are ordinary absolutely convergent integrals.  The last example
uses the short operator $(\mathsf E-\mathsf E^2)/2$ and
$\mathcal A_1(2)=-\pi^2/4$.

The results suggest several precise continuations.
\begin{enumerate}
\item Determine which noninteger beta weights retain a finite expression
in shifts of $\mathscr Z(s)$.  The present factored operator uses the
integrality of the denominator exponent, and does not settle finite
closure for a general complex exponent.
\item Extend the factored operator to kernels with several independent
negative poles.  Partial fractions give a finite reduction to scaled
instances of $\mathcal A_s$; an explicit confluent formula should keep
coalescing poles and order derivatives regular simultaneously.
\item Determine which bilinear full-line integrals with two independent
polylogarithm orders admit finite reductions of comparable simplicity.
A product introduces an additional convolution; the present one-factor
theorem does not settle that question.
\item Study the analytic continuation of the local series
\eqref{mpr:eq:localexpansion} around $u=\pm2\pi i$ and its monodromy.
The theorem proves normal convergence within the stated disk, but does
not supply a global single-valued continuation in $u$.
\item Compare the precise all-order transform and resonance formulas with
the older Mellin-transform tables and subsequent Euler-sum literature
before asserting priority beyond the inspected corpus.
\end{enumerate}

% END sections/04_mellin.tex


% BEGIN sections/05_stieltjes_edge.tex
\section{Completion of the undifferentiated Stieltjes edge}
\label{edge:sec}

The positive-argument-derivative product theorem in the earlier
\emph{Polygamma Tail Identities} report \cite{tail-polygamma} assumes that both ordinary
derivative orders are positive.  Its research question 9 asks for the
missing cases in which one factor is undifferentiated.  We resolve
that edge.  The mechanism remains the classical Mikol\'as product
integral; the additional result is the joint spectral-pole completion,
its all-index finite algorithm, and its exact endpoint law.  No claim
of historical priority for the classical product kernel is made.
The five incoming packages at the pinned October 11 revision were
also checked for duplication of this statement; their spectral
Stieltjes defects concern different products.

Use the convention
\[
 \zeta(1+u,x)=\frac1u+
       \sum_{m\ge0}\frac{(-1)^m}{m!}\gamma_m(x)u^m,
 \qquad \gamma_0(x)=-\psi(x).
\]
Every finite part below means the constant coefficient of the integral
from $\epsilon$ to $1$ in the original coordinate $x$.  Terms
$\epsilon^a\log^j\epsilon$ with $a<0$ and all positive powers of
$\log\epsilon$ are deleted.  A change of cutoff coordinate is not
part of this convention.

\subsection{The completed first-order edge}

For an integer $l\ge2$, put
\begin{align}
 P_l(u,v)&=\frac{\Gamma(l+u+v)}{\Gamma(1+u)\Gamma(l+v)}
                              \zeta(l+u+v),\\
 E(u)&=\frac{\pi u}{\sin\pi u},\\
 H_l(u,v)&=-\frac{P_l(u,v)E(u)-P_l(0,v)}u
       +(-1)^l\frac{P_l(u,v)E(v)-P_l(u,0)}v.
 \label{edge:H}
\end{align}
Each divided difference in \eqref{edge:H} is holomorphic near
$(0,0)$.  In particular $H_l$ is a genuine two-variable Taylor germ,
not an iterated choice of Laurent constants.

\begin{theorem}[Every undifferentiated-factor jet]
\label{edge:theorem}
For all integers $m,b\ge0$ and $l\ge2$,
\begin{equation}
 \left.\partial_u^m\partial_v^b H_l(u,v)\right|_{u=v=0}
 =(-1)^m\operatorname{FP}\int_0^1
    \gamma_m(x)\left.\partial_t^b\zeta(t,x)\right|_{t=l}\,dx.
 \label{edge:alljets}
\end{equation}
Consequently every finite part
\[
 \operatorname{FP}\int_0^1\gamma_m(x)\gamma_n^{(r)}(x)\,dx,
 \qquad m,n\ge0,\quad r\ge1,
\]
has a finite expression in ordinary zeta values and derivatives at
positive integers.  The superscript $(r)$ is an argument derivative.
\end{theorem}

\begin{proof}
Write $s=1+u$ and $t=l+v$.  The classical product integral continues as
\begin{align}
 F(s,t)&=\pi\frac{\Gamma(s+t-1)}{\Gamma(s)\Gamma(t)}
              \zeta(s+t-1)\{\csc(\pi s)+\csc(\pi t)\}.
 \label{edge:F}
\end{align}
Initially it equals $\int_0^1\zeta(s,x)\zeta(t,x)\,dx$ when
$\Re s<1$, $\Re t<1$, and $\Re(s+t)<1$; see
Shpot--Chaudhary--Paris, Remark 1 \cite{ShpotParis}, or prove it by integrating the
absolutely convergent Hurwitz Fourier series first for
$\Re s,\Re t<0$.  In that initial region the scalar identity
$\int_0^1\zeta(t,x)\,dx=0$ implies
\[
 \int_0^1\left(\zeta(s,x)-\frac1{s-1}\right)\zeta(t,x)\,dx
 =F(s,t).
\]
Both sides have the same meromorphic continuation.

Set $R(s,x)=\zeta(s,x)-1/(s-1)$.  For $s$ near $1$, this function
is holomorphic in $s$, and
\[
 R(s,x)=x^{-s}+A(s,x),\quad
 A(s,x)=\zeta(s,1+x)-\frac1{s-1}.
\]
The function $A$ is jointly holomorphic near $s=1$, $x=0$.
For the second factor use
\[
 \zeta(t,x)=x^{-t}+B(t,x),\quad B(t,x)=\zeta(t,1+x).
\]
In the product the two exponents crossing $-1$ at $(s,t)=(1,l)$ are
precisely
\[
 \zeta(t)x^{-s},\qquad
 c_{l-1}(s)x^{l-1-t},\quad
 c_j(s)=\frac{(-1)^j(s)_j}{j!}\zeta(s+j).
\]
The subtraction of $1/(s-1)$ changes only the constant Taylor
coefficient of $A$; since $l-1\ge1$, the second coefficient is still
$c_{l-1}(s)$ and is holomorphic near $s=1$.
After subtracting sufficiently many further Taylor monomials, the
remainder and all prescribed parameter derivatives are dominated by
$Cx^{-\delta}(1+|\log x|^N)$, with $\delta<1$.
Their integrals therefore form a holomorphic germ.

The analytic integrals of the two resonant monomials are
$-\zeta(t)/(s-1)$ and $-c_{l-1}(s)/(t-l)$, whereas all their
parameter derivatives at the resonance have coordinate finite part
zero: they are multiples of $x^{-1}\log^j x$.
Every nonresonant monomial commutes with parameter differentiation.
Thus the required holomorphic completion is
\[
 F(s,t)+\frac{\zeta(t)}{s-1}+
                         \frac{c_{l-1}(s)}{t-l}.
\]
Now $P_l(0,v)=\zeta(l+v)$ and
$c_{l-1}(1+u)=-(-1)^lP_l(u,0)$.
Substitution into \eqref{edge:F} gives exactly \eqref{edge:H}.
Finally $\partial_u^mR(1+u,x)|_{u=0}=(-1)^m\gamma_m(x)$,
proving \eqref{edge:alljets} with the stated finite-part convention.
\end{proof}

\subsection{A terminating coefficient algorithm}

Define
\[
 p_{r,j}=\left.\frac{d^j}{du^j}(1+u)_r\right|_{u=0},
 \qquad p_{r,j}=0\ \text{if }j>r.
\]
For $r\ge1$ the Hurwitz derivative law gives
\begin{equation}
 \gamma_n^{(r)}(x)=(-1)^{n+r}
       \sum_{j=0}^{\min(n,r)}\binom nj p_{r,j}
            \left.\partial_t^{n-j}\zeta(t,x)\right|_{t=r+1}.
 \label{edge:conversion}
\end{equation}
Write $P_l(u,v)=\sum_{a,b\ge0}P^{(l)}_{ab}u^av^b$ and
$E(u)=\sum_{j\ge0}e_{2j}u^{2j}$, where
\[
 e_0=1,\qquad e_{2j}=2(1-2^{1-2j})\zeta(2j)\quad(j\ge1).
\]
Then
\begin{equation}
 [u^av^b]H_l=
 -\sum_{j=0}^{\lfloor(a+1)/2\rfloor}e_{2j}P^{(l)}_{a+1-2j,b}
 +(-1)^l\sum_{j=0}^{\lfloor(b+1)/2\rfloor}
                        e_{2j}P^{(l)}_{a,b+1-2j}.
 \label{edge:finite}
\end{equation}
Combining \eqref{edge:conversion} and \eqref{edge:finite} evaluates
any requested $m,n,r$ with coefficients of total degree at most
$m+n+1$.  There are no spectral limits left in the algorithm.
The coefficient source can itself be written as
\begin{equation}
 P_l(u,v)=\zeta(l+u+v)
       \prod_{j=1}^{l-1}\left(1+\frac{u}{j+v}\right)
       \exp\left\{
       \sum_{k\ge2}\frac{(-1)^k\zeta(k)}k
                    [(u+v)^k-u^k-v^k]\right\}.
 \label{edge:coefficient-source}
\end{equation}
The product is finite, and only finitely many terms of the exponential
can affect a prescribed coefficient.  This exhibits the claimed
ordinary-zeta closure directly.
In conjunction with the inherited positive-positive sector this proves
a useful uniform statement: whenever the total ordinary derivative
order is positive, every quadratic same-point Stieltjes finite part
reduces to a finite expression involving ordinary positive-integer
zeta values and their order derivatives.  No Stieltjes constant
is needed in those final expressions.  The double-undifferentiated
corner is different and is recorded below for comparison.

For completeness the edge and inherited interior can be written in one
finite dictionary. For $r,t\ge0$ with $r+t>0$, put
$k=r+1$, $l=t+1$, $w=r+t+1$, and define
\begin{align}
 P_{kl}(u,v)&=\frac{\Gamma(w+u+v)}{\Gamma(k+u)\Gamma(l+v)}
                         \zeta(w+u+v),\notag\\
 H_{kl}(u,v)&=(-1)^k\frac{P_{kl}(u,v)E(u)-P_{kl}(0,v)}{u}
       +(-1)^l\frac{P_{kl}(u,v)E(v)-P_{kl}(u,0)}{v}.
       \label{edge:uniform-kernel}
\end{align}
Then the uniform coefficient formula is
\begin{equation}\label{edge:uniform-dictionary}
 \boxed{\displaystyle
 \operatorname{FP}\int_0^1\gamma_m^{(r)}(x)\gamma_n^{(t)}(x)\,dx
 =(-1)^{m+n+r+t}m!n![u^mv^n]
             (1+u)_r(1+v)_tH_{kl}(u,v).}
\end{equation}
When $r=0$ this is the new edge theorem; when $t=0$ use symmetry.
When both are positive it is the inherited positive-positive formula
\cite[Theorem \texttt{hprod:stieltjes-products}]{tail-polygamma}.
Multiplying its Hurwitz Taylor germ by the two rising factorials
proves \eqref{edge:uniform-dictionary}. Only order derivatives at the
single argument $w=r+t+1$ occur. Their coefficients involve rational
numbers and ordinary zeta values from the Gamma expansion. Thus
``single argument'' refers to the differentiated zeta factors, not
to every undifferentiated zeta value in the coefficients.

\subsection{Explicit first-index identities}

Let $H=H_{l-1}$, $H_2=H_{l-1}^{(2)}$, and
$\sigma=(-1)^l$.  The first coefficients are
\begin{align}
 H_l(0,0)&=-H\zeta(l)+(\sigma-1)\zeta'(l),\\
 \partial_u H_l(0,0)
 &=\left[(\sigma-1)\zeta(2)+\left(\frac12-\sigma\right)H_2
                          -\frac12H^2\right]\zeta(l)
       +(\sigma-1)H\zeta'(l)
       +\left(\sigma-\frac12\right)\zeta''(l),
       \label{edge:firstu}\\
 \partial_v H_l(0,0)
 &=\{(\sigma-1)\zeta(2)+H_2\}\zeta(l)
       -H\zeta'(l)+\left(\frac\sigma2-1\right)\zeta''(l).
       \label{edge:firstv}
\end{align}
For example the $l=2$ formulas give the pair
\begin{align}
 \boxed{\operatorname{FP}\int_0^1
       \gamma_1(x)\gamma_0'(x)\,dx
          =\frac12\zeta''(2)-\zeta(2),}
       \label{edge:example10}\\
 \boxed{\operatorname{FP}\int_0^1
       \gamma_0(x)\gamma_1'(x)\,dx
          =-\zeta'(2)-\frac12\zeta''(2).}
       \label{edge:example01}
\end{align}
Neither individual identity follows merely by integrating the
derivative of $\gamma_0\gamma_1$.  Their sum is the endpoint
constant $-\zeta(2)-\zeta'(2)$.
At one further argument derivative,
\begin{equation}
 \operatorname{FP}\int_0^1\gamma_1(x)\gamma_0''(x)\,dx
  =\left(4\zeta(2)-\frac32\right)\zeta(3)
          +6\zeta'(3)+3\zeta''(3).
 \label{edge:example12}
\end{equation}

One higher Stieltjes index already produces a new derivative order:
\begin{equation}
 \operatorname{FP}\int_0^1\gamma_2(x)\gamma_0'(x)\,dx
 =2\zeta'(2)-\zeta''(2)-\frac23\zeta'''(2)
                         +2\zeta(2)\zeta(3).
 \label{edge:example20}
\end{equation}
More generally, put $J_{0,1}^{m,n}=\operatorname{FP}\int_0^1
\gamma_m(x)\gamma_n'(x)\,dx$.  In its finite formula the
coefficient of the highest ordinary-zeta
order derivative $\zeta^{(m+n+1)}(2)$ is exactly
\begin{equation}
 \frac{(-1)^{m+n+1}(m-n)}{(m+1)(n+1)}.
 \label{edge:top}
\end{equation}
This is a coefficient statement for the displayed reduction and does
not assert arithmetic independence of its evaluated constants.
Indeed the highest derivative can only come from
$-P^{(2)}_{m+1,n}+P^{(2)}_{m,n+1}$ with all Gamma and
Pochhammer factors replaced by their constant terms.  Multiplication
by $(-1)^{m+n+1}m!n!$ gives \eqref{edge:top}.

For readers preferring an ordinary convergent integral, put
$g(x)=\gamma_1(1+x)$ and $h(x)=\psi'(1+x)$.
Equation \eqref{edge:example10} is equivalent to
\begin{equation}
 \begin{split}
 \int_0^1\left\{
   \frac{g(x)-\gamma_1-(\zeta(2)+\zeta'(2))x}{x^2}
  +\frac{\log x}{x}\{h(x)-\zeta(2)\}+g(x)h(x)
                  \right\}\,dx
 =\zeta(2)-\frac12\zeta''(2)+\frac14+\gamma_1.
 \end{split}
 \label{edge:ordinary}
\end{equation}
The first apparent singularity is removable and the second is
$O(\log x)$.  Indeed
$\gamma_1(x)=\log x/x+g(x)$ and $\psi'(x)=x^{-2}+h(x)$;
the removed singular monomials have finite-part integral
$-1/4-\gamma_1$, since those proportional to $x^{-1}$ or
$x^{-1}\log x$ have finite part zero.  This proves
\eqref{edge:ordinary} without any numerical limiting prescription.

\subsection{All endpoint constants and compatibility with integration by parts}

For $r,t,m,n\ge0$ write
\[
 J_{r,t}^{m,n}=\operatorname{FP}\int_0^1
                         \gamma_m^{(r)}(x)\gamma_n^{(t)}(x)\,dx.
\]
Set $p_{0,0}=1$ and $p_{0,j}=0$ for $j>0$, consistently with
$(1+u)_0=1$.

\begin{theorem}[Exact endpoint law]
\label{edge:endpoint}
For every $r,t,m,n\ge0$,
\begin{equation}
 \boxed{\begin{aligned}
 J_{r+1,t}^{m,n}+J_{r,t+1}^{m,n}
 ={}&-\frac{(-1)^{m+r}p_{r,m}}{(r+1)!}
                       \gamma_n^{(r+t+1)}(1)\\
    &-\frac{(-1)^{n+t}p_{t,n}}{(t+1)!}
                       \gamma_m^{(r+t+1)}(1).
 \end{aligned}}
 \label{edge:ibp}
\end{equation}
All constants on the right reduce to positive-integer ordinary-zeta
jets by \eqref{edge:conversion}.
\end{theorem}

\begin{proof}
The shift identity gives
$\gamma_m(x)=x^{-1}\log^m x+\gamma_m(1+x)$ exactly.
Its $r$th derivative has singular part
\[
 S_{m,r}(x)=\frac{(-1)^r}{x^{r+1}}
      \sum_{j=0}^{\min(m,r)}(-1)^j\binom mjp_{r,j}\log^{m-j}x.
\]
The coefficient of $x^{-r-1}\log^0x$ is
$(-1)^{m+r}p_{r,m}$; it vanishes if $m>r$.
The smooth factor is $\gamma_m^{(r)}(1+x)$.
The constant term of the product at zero has three contributions:
the product of its two smooth constant terms, and the two products
of a singular logarithm-free coefficient with the matching smooth
Taylor coefficient.  The singular--singular product has power
$x^{-r-t-2}$ and contributes no constant.
The first contribution cancels the endpoint value at $x=1$.
The remaining two are exactly the negative of the right side of
\eqref{edge:ibp}.  Apply the fundamental theorem of calculus on
$[\epsilon,1]$ and take constant coefficients.
\end{proof}

This law provides a useful consistency check on the completed kernels.
It also shows precisely why the first-derivative edge contains information
beyond the old positive-positive sector: integration by parts determines
$J_{0,1}^{m,n}+J_{1,0}^{m,n}$, but not its antisymmetric part.  Once
the first edge has been established, higher edges can also be related
to the positive-positive formulas while retaining the explicit endpoint
terms.  In particular no formal ``integral of a derivative is zero''
rule is appropriate here.

\subsection{The double-undifferentiated corner}

The corner $r=t=0$ is already treated in the manuscript's
coincident-endpoint Stieltjes chapter.  For compatibility, a useful
holomorphic presentation is as follows.  Define
\begin{align}
 A(u)&=u\zeta(1+u),\\
 C(u,v)&=\frac{\Gamma(1+u+v)}{\Gamma(1+u)\Gamma(1+v)}
   \frac{\pi uv\{\csc(\pi u)+\csc(\pi v)\}}{u+v}
                                    A(u+v).
\end{align}
All apparent singularities in $C$ are removable near zero.  In
particular $C(u,0)=A(u)$ and $C(0,v)=A(v)$.
The completed corner is
\begin{equation}
 H_{11}(u,v)=\frac{A(u)+A(v)-1-C(u,v)}{uv}.
 \label{edge:corner}
\end{equation}
Its Taylor coefficients satisfy
$J_{0,0}^{m,n}=(-1)^{m+n}m!n![u^mv^n]H_{11}$.
This expression follows from the same local subtraction, now applied
to both factors $\zeta(1+u,x)-1/u$ and
$\zeta(1+v,x)-1/v$: the integrated spectral product is
$F(1+u,1+v)+1/(uv)$, and the local resonance corrections are
$(\zeta(1+v)-1/v)/u$ and $(\zeta(1+u)-1/u)/v$.
The two pole subtractions must be made jointly.
Equation \eqref{edge:corner} records compatibility with the inherited
corner rather than claiming it as another new result.


\subsection{Independent verification}

The supplied script constructs the local Gamma germ with exact
rational arithmetic and formal ordinary-zeta derivatives.  It checks
all 81 endpoint identities with $0\le m,n,r,t\le2$, compares the
printed low-index identities, and exports their formulas.  This is
an exact finite verification of the implementation, not a replacement
for the analytic proof at all indices.
The numerical suite computes the left sides by integrating their
convergent local Laurent--logarithmic expansions up to a positive
splitting point and then using ordinary quadrature to $1$.
Thus it does not evaluate the product kernel to obtain both sides
of a test.  Its precision, truncation order, actual residuals, and
second splitting-point check are recorded in the JSON output.
These arbitrary-precision comparisons are diagnostics rather than
interval certificates.

% END sections/05_stieltjes_edge.tex


% BEGIN sections/06_order_ten.tex
% Research fragment. Source statements checked against the incoming files
% Alternating_Tails_Spectral_Defects and Quartic_Spectral_Reconstruction.
% Suggested bibliography keys: tail-alternating, tail-quartic.
\section{Complete optimal reconstruction at smaller tail order ten}
\label{t10:section}

Two recent source theorems leave only one pair at smaller order ten.
The order-ten theorem of \emph{Alternating Tails and Spectral Defects},
Theorem~\texttt{rt:rank10}, proves the conjecture for every integer
$b\ge6$ outside the progression $b\equiv962\pmod{2048}$.
The effective theorem of \emph{Quartic Spectral Reconstruction},
Theorem~\texttt{ert:main} and its exact threshold calculation,
proves it for every $b\ge1920$. We certify the single remaining
value $b=962$ by a five-square determinant over a finite field.
This completes the whole smaller-order-ten family. The method does
not assert that the unrestricted unequal-order conjecture is proved.

\subsection{The formal coordinate problem}

For even $p,q$ put
\[
 A_p(n)=\zeta(p,n+1)=\zeta(p)-H_n^{(p)},\qquad x_n=n+\tfrac12,
\]
and initially define
\[
 D_{p,q}(s)=\sum_{n\ge0}A_p(n)A_q(n)x_n^{-s},
 \qquad \Re s>3-p-q.
\]
Let $D_r=D_{p,q}(-2r)$ denote its regular continued centered moments.
The centered Hurwitz expansion has powers $x_n^{1-p-2j}$; hence
the product has only even inverse powers. Its possible continuation
poles lie at odd integers, so all the displayed $-2r$ values are regular.
The inherited centered continuation and finite Bernoulli reduction
supply inverse transforms
\begin{align}
 C_j&=\sum_{r=0}^j\binom{2j+1}{2r}2^{-2(j-r)}D_r,
 \label{t10:inverse}\\
 D_r&=\sum_{j=0}^r\frac{\binom{2r}{2j}}{2j+1}
 B_{2r-2j}(1/2)C_j.\label{t10:forward}
\end{align}
Write $p=2a$, $q=2b$, $a<b$. To specify the algebraic claim completely,
let $\mathcal V_{a,b}$ be the $2b$-dimensional rational vector space with
\emph{formal basis}
\[
 C_0,\ldots,C_{b-1},Z_0,\ldots,Z_{b-1}.
\]
These are independent labels by definition. Under numerical evaluation,
$Z_k$ is assigned $\zeta(2k)$, including $Z_0=-1/2$, and each early
$C_j$ is assigned its continued value above. No injectivity of this
evaluation map is asserted. For every $j\ge b$, the source reduction
defines $C_j=\sum_{k=0}^{b-1}M_{j-b,k}Z_k$, using the formula below
also for row indices greater than $b-1$. Thus every later $D_r$ is a
specified vector of $\mathcal V_{a,b}$. The coordinate matrix of
$C_0,\ldots,C_{2b-1}$ is $\operatorname{diag}(I_b,M_{a,b})$.
It follows that the first $2b$ moment coordinates have full formal
rank precisely when the matrix
\begin{equation}\label{t10:M}
 (M_{a,b})_{ik}=B_{2(i+k-a+1)}
 \left\{\frac{\binom{2i+2}{2a-2k}}{2i+2}
 +\frac{\binom{2b+2i-2a+2}{2b-2k}}{2b+2i-2a+2}\right\},
 \qquad0\le i,k<b,
\end{equation}
is nonsingular. A negative Bernoulli index gives zero. A binomial
summand is omitted unless its lower index is at least two and no
greater than its upper index. Thus a lower index zero is omitted.
The complete first-$2b$ coordinate determinant is
$\det M_{a,b}/(4b-1)!!$. These are inherited identities, not new
independence assertions about numerical zeta values.

\begin{theorem}[Complete smaller-order-ten reconstruction]
\label{t10:complete}
For every integer $b\ge6$, $M_{5,b}$ is nonsingular. Consequently,
for every even $q>10$, the moments
\[
 D_{10,q}(0),D_{10,q}(-2),\ldots,D_{10,q}(-2q+2)
\]
have formal coordinate rank $q$ in the inherited coordinate space.
Every later centered moment is a finite rational linear combination
of these first $q$ moments.
\end{theorem}

The only new nonvanishing certificate needed by the theorem is
\begin{equation}\label{t10:main-certificate}
 \boxed{\det M_{5,962}\equiv15578\pmod{65537}.}
\end{equation}
The rational determinant has denominator prime to $65537$;
therefore this congruence proves exact nonvanishing.

\subsection{The retained five-square matrix}

We state the inherited contraction explicitly to make the finite
certificate independently reproducible. For $b\ge2a$ and
$0\le i,\ell<a$, define
\[
 n_i=b-a+i,\quad N_{i\ell}=2(n_i-\ell)+1,\quad
 R_\ell=2b-2a+2+2\ell,
\]
\[
 F_{i\ell}=\frac1{(2\ell+1)!}
 \prod_{r=0}^{2\ell-1}(2n_i-r)(2n_i+3+r).
\]
Define the integer Newton coefficients $d_{ih}$ by
\begin{equation}\label{t10:Newton}
 (2n_i+2-y)(4n_i+2-y)^{\underline{2i}}
       =\sum_{h=0}^{2i+1}d_{ih}y^{\underline h}.
\end{equation}
Here $z^{\underline h}=z(z-1)\cdots(z-h+1)$. Equivalently,
if the polynomial on the left is $W_i(y)$, then
\[
 d_{ih}=\frac1{h!}\sum_{j=0}^h(-1)^{h-j}\binom hj W_i(j).
\]
The retained matrix is
\begin{equation}\label{t10:A}
 \mathcal A_{i\ell}=2F_{i\ell}
 +\frac{R_\ell!}
 {(2n_i+1)(2n_i+2)(2\ell+2)!N_{i\ell}!}
 \sum_{k=0}^i d_{i,2k+1}N_{i\ell}^{\underline{2k+1}}
                 2B_{2b-2(a-i+\ell+k)}.
\end{equation}
The omitted lower-index-zero convention in \eqref{t10:M} is
already included in the $2F$ term and must not be altered.

For completeness, the source multiplier is
\begin{equation}\label{t10:multiplier}
 \begin{aligned}
 \det M_{a,b}={}&(-1)^{a(b-a)+a(a-1)/2}
 \frac{\prod_{t=0}^{b-1}(2b+2t-2a+1)!}
      {\prod_{k=0}^{b-1}(2b-2k)!}\\
 &\quad\cdot\prod_{n=0}^{b-a-1}\frac{((2n+1)!)^2}{(4n+2)!(4n+3)!}
 \prod_{i=0}^{a-1}\frac{(2n_i+1)(2n_i+2)!}{(4n_i+2)!}
 \det\mathcal A.
 \end{aligned}
\end{equation}
The exact derivation is the Pad\'e-row elimination in the source's
Lemma~\texttt{ert:normalization}. It can also be replayed directly:
the monic rows
\[
 p_{nt}=\frac{2(2n+2t+1)!(2n+1)!}
 {(4n+2)!(2n-2t+1)!(2t)!},\qquad0\le t\le n,
\]
are applied to the factorial-normalized Bernoulli matrix. Their
orthogonality eliminates the last $b-a$ columns from the final $a$
rows. Reversing the retained columns, and dividing row $i$ by
$c_{n_i}=(2n_i+1)(2n_i+2)!/(4n_i+2)!$, gives \eqref{t10:A}.
The diagonal eliminated entries are
$((2n+1)!)^2/((4n+2)!(4n+3)!)$, which gives
\eqref{t10:multiplier}, including its sign.

\subsection{The complete finite certificate}

Set $a=5$, $b=962$, and $P=65537$. The prime $P$ exceeds every
factorial argument in \eqref{t10:A} and \eqref{t10:multiplier}.
The Bernoulli recurrence
\begin{equation}\label{t10:Brec}
 B_0=1,\qquad
 B_n=-\frac1{n+1}\sum_{j=0}^{n-1}\binom{n+1}{j}B_j
 \quad(n\ge1)
\end{equation}
shows that every $B_n$ used here has denominator prime to $P$:
all indices are at most $1922<P-1$. Thus the entire calculation
can be performed in $\mathbb F_P$ without an undefined division.
Using $B_1=-1/2$ and $B_{2r+1}=0$ for $r\ge1$, the nine required
residues, in descending-index order, are
\begin{equation}\label{t10:Bvalues}
\begin{array}{c|rrrrrrrrr}
2r&1922&1920&1918&1916&1914&1912&1910&1908&1906\\\hline
 B_{2r}\bmod P&31528&45962&26974&45003&39580&2844&23974&18828&26318.
\end{array}
\end{equation}
Substitution into \eqref{t10:A} gives
\begin{equation}\label{t10:Amatrix}
\mathcal A\equiv
\begin{pmatrix}
25959&60986&10404&1273&2935\\
16142&62424&8365&2781&39264\\
13624&13203&159&63417&38163\\
51637&50007&47664&40198&52007\\
23028&22542&12445&19571&47569
\end{pmatrix}\pmod{65537}.
\end{equation}
Gaussian elimination without row exchanges has diagonal pivots
\[
25959,\quad2542,\quad6997,\quad61613,\quad26239.
\]
Their product is $60407$ modulo $65537$. The complete multiplier
in \eqref{t10:multiplier}, including its sign, is $49105$ modulo
$65537$. Therefore
\[
\det M_{5,962}\equiv49105\cdot60407\equiv15578\not\equiv0
     \pmod{65537},
\]
which proves \eqref{t10:main-certificate}.

\begin{proof}[Proof of Theorem~\ref{t10:complete}]
If $b\not\equiv962\pmod{2048}$, use the source order-ten theorem.
If $b\ge1920$, use the source effective eventual theorem with
$B_5=1920$. The only integer $b$ with
$6\le b<1920$ and $b\equiv962\pmod{2048}$ is $962$;
\eqref{t10:main-certificate} treats that value. These cases cover
all integers $b\ge6$. The inherited triangular transforms and
finite late-coordinate expansion then give the stated moment
reconstruction.\end{proof}

The supplied verifier \texttt{complete\_p10.py} implements
\eqref{t10:Brec}, \eqref{t10:Newton}, \eqref{t10:A}, and modular
Gaussian elimination, using only Python integer arithmetic. A
second verifier reconstructs the same matrix from the uncontracted
Pad\'e row sums and separately reduces exact rational Bernoulli
numbers obtained from SymPy. These checks agree entry by entry.
The prime is checked by trial division through its integer square
root. No floating-point determinant is used anywhere in the proof.
This is a finite arithmetic proof certificate, not a claim of
proof-assistant formalization.

\subsection{An ordinary convergent polygamma identity for every larger order}

The nonvanishing theorem makes the inherited reconstruction an
unconditional all-$q$ identity family at smaller order ten. Define,
for $0\le k<b$,
\begin{equation}\label{t10:weight}
 P_{b,k}(x)=\sum_{t=0}^{b-1}(M_{5,b}^{-1})_{kt}
 \left\{(x+\tfrac12)^{2b+2t+1}-(x-\tfrac12)^{2b+2t+1}\right\}
 =\sum_{r=0}^{2b-1}w_{b,k,r}x^{2r}.
\end{equation}
All coefficients are rational. Put
\[
 s_{p,0}=\frac1{p-1},\qquad
 s_{p,j}=\frac{B_{2j}(1/2)}{(2j)!}(p)_{2j-1}\quad(j\ge1),
 \qquad
 \alpha_j=\sum_{h=0}^js_{10,h}s_{2b,j-h},
\]
where $(p)_m$ is the rising factorial. Define the finite even polynomial
\begin{equation}\label{t10:subtraction}
 S_{b,k}(x)=\sum_{r=b+4}^{2b-1}w_{b,k,r}
          \sum_{j=0}^{r-b-4}\alpha_jx^{2(r-b-4-j)}.
\end{equation}
Then
\begin{equation}\label{t10:ordinary}
 \boxed{\sum_{n=0}^{\infty}
 \left\{P_{b,k}(x_n)\,
 \frac{\psi^{(9)}(n+1)\psi^{(2b-1)}(n+1)}{9!\,(2b-1)!}
                -S_{b,k}(x_n)\right\}=\zeta(2k).}
\end{equation}
Every sum in \eqref{t10:ordinary} converges absolutely. Indeed,
the centered Hurwitz expansion is
\[
 A_p(n)\sim\sum_{j\ge0}s_{p,j}x_n^{1-p-2j}.
\]
Thus \eqref{t10:subtraction} removes exactly every nondecaying
term of $P_{b,k}(x_n)A_{10}(n)A_{2b}(n)$ and leaves $O(x_n^{-2})$.
The polynomial difference in \eqref{t10:weight} is exactly the
inverse transform \eqref{t10:inverse}. Matrix inversion therefore
makes its continued sum $\zeta(2k)$. All removed monomials have
zero centered continued sum, since
$\zeta(-2j,1/2)=0$ for integers $j\ge0$, proving the ordinary
convergent identity. The polygamma substitution uses
$\psi^{(p-1)}(n+1)=(p-1)!\zeta(p,n+1)$ for even $p$.

The new contribution is the complete smaller-order-ten theorem,
including its single missing exact certificate. The centered
continuation, Pad\'e contraction, and conversion of reconstructed
coordinates into convergent sums retain their source attribution.

% END sections/06_order_ten.tex


% BEGIN sections/07_order_twelve.tex
% Research fragment. Requires the source-coordinate definitions in
% order_ten_completion.tex, or equivalent definitions from the parent article.
% No external computational package is required by the complete replay.
\section{Complete optimal reconstruction at smaller tail order twelve}
\label{t12:section}

Research question R1 of the source \emph{Quartic Spectral Reconstruction}
explicitly asks for complete smaller-order-ten and smaller-order-twelve
families. The preceding section and the theorem below resolve those two
specified tasks. At smaller order twelve, the preceding one-case synthesis
is unavailable:
no source theorem handles almost every parameter. We instead prove an
explicit eventual threshold by an elementary coefficient estimate, then
certify the entire remaining finite set over one finite field. The proof
uses only six-square retained matrices, even when the original matrix has
more than eleven thousand rows.

\begin{theorem}[Complete smaller-order-twelve reconstruction]
\label{t12:complete}
For every integer $b\ge7$, the matrix $M_{6,b}$ defined in
\eqref{t10:M} is nonsingular. Consequently, for every even $q>12$,
the first $q$ centered moments of $A_{12}(n)A_q(n)$ have formal
coordinate rank $q$, and every later centered moment is a finite
rational linear combination of them.
\end{theorem}

The proof has two exact components:
\[
 \det M_{6,b}\ne0\quad(b\ge11520),\qquad
 \det M_{6,b}\ne0\quad(7\le b\le11519).
\]
The first component has an ordinary analytic proof below. The second is
a finite arithmetic certificate with $11513$ cases, reproduced in full
by two independently implemented verifiers. No sampling of the parameter
range is used.

\subsection{A six-square polynomial and its leading determinant}

In \eqref{t10:Newton}--\eqref{t10:A}, set $a=6$ and replace
$2B_{2b-2j}$ by independent indeterminates $U_j$, $1\le j\le11$.
Write the resulting six-square matrix as $\mathcal A(b;U)$. Set
\begin{align*}
 n_i&=b-6+i,&R_\ell&=2b-10+2\ell,\\
 J(b)&=\prod_{i=0}^5(2n_i+1)(2n_i+2),&
 E(b)&=\prod_{i,\ell=0}^5
       \prod_{r=1}^{\max(0,2i-4\ell-1)}(R_\ell+r),\\
 D(b)&=12!J(b)E(b),&
 \mathscr P(b;U)&=\det(D(b)\mathcal A(b;U)).
\end{align*}
The polynomial $E$ has $35$ linear factors. The source factorial
cancellation proves
\[
 D\mathcal A\in\mathbb Z[b,U_1,\ldots,U_{11}]^{6\times6},
 \qquad\deg_b(D\mathcal A_{i\ell})\le67,
 \qquad\deg_b\mathscr P\le402.
\]
Indeed $N_{i\ell}-R_\ell=2i-4\ell-1$, so $E$ clears exactly the
possible denominator products in $R_\ell!/N_{i\ell}!$;
$J$ clears the row factors and $12!$ clears the fixed factorials.
All these factors are positive when $b\ge12$.

Let $\mathscr P_{[6]}$ denote the part homogeneous of degree six in
the $U$ variables. The inherited highest-degree determinant identity,
specialized to six, gives at $U_j=2B_{2b-2j}$
\begin{equation}\label{t12:leading}
 \left|\mathscr P_{[6]}\right|
 =\frac{\mathcal G(b)(4F_0)^6}{(2\pi)^{12(b-6)}}
   \det\bigl(\zeta(s_0+2i+2j)\bigr)_{0\le i,j<6},
 \qquad s_0=2b-22,\quad F_0=s_0!,
\end{equation}
where
\begin{align*}
 \mathcal G(b)={}&\kappa J(b)^5E(b)^6
 \prod_{i=0}^5(4n_i+2)^{\underline{2i}}
 \prod_{\ell=0}^5\prod_{r=1}^{12+2\ell}(2b-22+r),\\
 \kappa={}&\frac{6!(12!)^6}{\prod_{\ell=0}^5(2\ell+2)!}
 =3590542128270039916216320000000.
\end{align*}
The degree of $\mathcal G$ is $402$, and every linear factor displayed
in it is at least $b$ when $b\ge24$. Hence
\begin{equation}\label{t12:Gbound}
 \mathcal G(b)\ge\kappa b^{402}\qquad(b\ge24).
\end{equation}

Here the leading determinant is a structural identity, not an estimate
of an expanded polynomial. Its short derivation uses
$N_{i\ell}-(2k+1)=2b-2(6-i+\ell+k)$ in \eqref{t10:A} to cancel
factorials. The resulting linear-in-$U$ part factors through the
matrix $(V_{6-i+\ell})$, where
$V_j=U_jF_0/(2b-2j)!$. The remaining row transformation is triangular,
with diagonal
$d_{i1}=-(i+1)(4n_i+2)^{\underline{2i}}$.
Taking determinants, reversing rows, and applying Euler's formula
$2B_{2m}=(-1)^{m-1}4(2m)!\zeta(2m)/(2\pi)^{2m}$ gives
\eqref{t12:leading} with the displayed $\mathcal G$. This retains the
attribution of the general highest-degree factorization in the source
\emph{Quartic Spectral Reconstruction}.

The zeta determinant in \eqref{t12:leading} is a positive moment
Gram determinant: use the points $r^{-2}$ with weights $r^{-s_0}$,
$r\ge1$. Cauchy--Binet and the first six distinct support points give
\begin{equation}\label{t12:zeta-lower}
 \begin{aligned}
 \det\bigl(\zeta(s_0+2i+2j)\bigr)_{0\le i,j<6}
 &\ge V\,720^{-s_0},\\
 V&=\prod_{1\le r<s\le6}(r^{-2}-s^{-2})^2\\
 &=\frac{14235529}{3743906242624487424000000000000}>0.
 \end{aligned}
\end{equation}
In particular, the highest-degree component cannot vanish.

\subsection{A simple coefficient estimate requiring no expanded determinant}

For an integer polynomial let $\|\cdot\|_1$ be the sum of absolute
values of its coefficients. It is submultiplicative. We shall use
the deliberately generous bound
\begin{equation}\label{t12:norm}
 \max_{0\le i,\ell<6}\|D\mathcal A_{i\ell}\|_1<2^{512}.
\end{equation}
The following elementary factor count proves it.
Every factor in $J$ has norm less than $16$, and every factor in $E$
has norm less than $16$. Thus
\[
 \|J\|_1<2^{48},\qquad\|E\|_1<2^{140},\qquad12!<2^{29}.
\]
After removing the two factors for one fixed row, the remaining
$J$ product has norm less than $2^{40}$. After cancellation in
$ER_\ell!/N_{i\ell}!$, at most $21$ additional linear factors are
present, each of norm less than $32$. Its norm is therefore less
than $2^{245}$. Also
\[
 \|N_{i\ell}^{\underline h}\|_1<2^{55}
 \quad(1\le h\le11),\qquad
 \|d_{ih}\|_1<2^{65}.
\]
For the last assertion, apply the Newton formula to
\[
 W_i(j)=(2b-10+2i-j)
             \prod_{r=0}^{2i-1}(4b-22+4i-j-r),
 \qquad0\le j\le h\le2i+1\le11.
\]
Its first factor has norm below $16$ and each remaining factor
has norm below $32$, so $\|W_i(j)\|_1<2^{54}$. The sum of the
binomial coefficients in the Newton formula is $2^h\le2^{11}$;
division by $h!$ only decreases the bound.

The numerator defining $F_{i\ell}$ has at most $20$ linear factors,
each of norm below $32$. Thus the constant term $2DF_{i\ell}$ has
norm below $2^{318}$. Each coefficient of an indeterminate $U_j$
in $D\mathcal A_{i\ell}$ has norm below
\[
 2^{29+40+245+65+55}=2^{434}.
\]
There are at most six such coefficients. Consequently the entry norm
is below $2^{318}+6\cdot2^{434}<2^{437}<2^{512}$, proving
\eqref{t12:norm}. This argument also proves integrality by the same
explicit cancellations; no fitted coefficient bound is involved.
A supplementary exact script also computes a sharper factor-product
bound; the proof uses only the elementary bound just established.

The determinant coefficient norm is at most
\[
 K=6!(2^{512})^6=720\,2^{3072}.
\]
For $b\ge11520$ put
\[
 T=\frac{8(2b)^{20}F_0}{(2\pi)^{s_0}}>1.
\]
The inequality $T>1$ follows from $s_0!/(2\pi)^{s_0}>(s_0/24)^{s_0}>1$,
using $s_0!>(s_0/3)^{s_0}$ and $2\pi<8$.
Euler's Bernoulli formula and $\zeta(2m)<2$ show
$|2B_{2b-2j}|\le T$ for $1\le j\le11$.
Since all remaining terms of $\mathscr P$ have total $U$ degree
at most five,
\[
 |\mathscr P_{[\le5]}|\le Kb^{402}T^5.
\]
Combining this estimate with \eqref{t12:leading}--\eqref{t12:zeta-lower},
then using $2\pi<8$, gives the explicit domination ratio
\begin{equation}\label{t12:ratio}
 \frac{|\mathscr P_{[\le5]}|}{|\mathscr P_{[6]}|}
 \le C\,\frac{b^{100}5760^{2b}}{(2b-22)!},\qquad
 C=\frac{2^{3289}}{\kappa V\,720^{21}}<2^{3289}.
\end{equation}
The final inequality follows, for example, from the exact rational value
\[
 \kappa V=\frac{786377640075797}{57600000}>1.
\]

\subsection{A threshold of 11520 proved by elementary inequalities}

Define
\[
 R(b)=2^{3289}\frac{b^{100}5760^{2b}}{(2b-22)!}.
\]
We prove $R(b)<1$ for every integer $b\ge11520$. At $b=11520$,
write $n=23018$. The elementary factorial estimate
$n!\ge(n/e)^n>(n/3)^n$, obtained by integrating $\log x$, gives
\[
 n!>7200^n=5760^n(5/4)^n,
\]
since $n>3\cdot7200$. Moreover $11520<2^{14}$,
$5760<2^{13}$, and $(5/4)^4=625/256>2$. Therefore
\begin{equation}\label{t12:initial-bound}
 R(11520)
 <2^{3289+1400+286}(4/5)^{23018}
 <2^{4975-5754}=2^{-779}<1.
\end{equation}
For every subsequent integer $b$,
\[
 \frac{R(b+1)}{R(b)}
 =\left(1+\frac1b\right)^{100}
       \frac{5760^2}{(2b-20)(2b-21)}<\frac34,
\]
because $b\ge100$ gives the first factor less than $3$,
both denominator factors exceed $b$, and $b\ge11520=2\cdot5760$.
Thus the ratio in \eqref{t12:ratio} is strictly below one for all
$b\ge11520$. Its highest-degree component cannot be cancelled.
As $D(b)>0$ and the source multiplier \eqref{t10:multiplier} is
nonzero, this proves $\det M_{6,b}\ne0$ throughout that infinite
range. The explicit threshold follows from a mathematical argument,
not from the observed behavior of a finite list of determinants.

\subsection{The complete finite-field certificate below the threshold}

For the remaining $11513$ integers $7\le b\le11519$, use the
single prime $P=65537$. The original source matrix is used for
$7\le b\le11$, where the contraction is not asserted. For
$12\le b\le11519$, use $\mathcal A$ from \eqref{t10:A} with $a=6$.
The source relation \eqref{t10:multiplier} implies that nonvanishing
of this six-square determinant suffices.

Every division is defined modulo $P$. The largest Bernoulli index
needed for the retained matrices is $2\cdot11519-2=23036<P-1$;
\eqref{t10:Brec} therefore puts all these Bernoulli numbers in
$\mathbb Z_{(P)}$. Every factorial argument and every positive
linear denominator appearing in the matrices and the exact source
multipliers is less than $4\cdot11520+5=46085<P$. In particular,
the entire source multiplier is a unit modulo $P$. The computation
never replaces a singular reduction modulo $P$ by an unsupported
claim of rational nonvanishing: every recorded residue is nonzero.

The file \texttt{complete\_p12.json} records the full ordered residue
list, with its first entry corresponding to $b=7$. It has $11513$
entries, and every entry belongs to $\{1,\ldots,65536\}$. For orientation,
a few entries are
\begin{equation}\label{t12:sample}
\begin{array}{c|rrrrrrrr}
 b&7&8&11&12&100&962&4096&11519\\\hline
\text{residue}&39450&1561&25713&60062&43034&42605&25983&23212.
\end{array}
\end{equation}
For $b\le11$ the entry is $\det M_{6,b}\bmod P$;
for $b\ge12$ it is $\det\mathcal A\bmod P$.
The displayed sample is explanatory; the proof checks the complete
ordered list.

Two full implementations replay the certificate by different exact
algorithms. The C++ verifier uses the triangular Bernoulli recurrence,
finite differences for the Newton coefficients at each integer $b$,
and modular Gaussian elimination. The Python verifier instead
inverts $(e^t-1)/t$ as a truncated power series, using exact integer
Kronecker multiplication. This produces $B_n/n!$ directly. It then
uses the factorial-cancelled entry formula
\begin{equation}\label{t12:cancelled-entry}
 \mathcal A_{i\ell}=2F_{i\ell}
 +\frac{R_\ell!}{(2n_i+1)(2n_i+2)(2\ell+2)!}
 \sum_{k=0}^i d_{i,2k+1}
 \frac{2B_{2b-2(6-i+\ell+k)}}{(2b-2(6-i+\ell+k))!},
\end{equation}
with Newton coefficients computed as global integer polynomials
in $b$, and evaluates determinants by a division-free subset
Laplace expansion. It verifies every one of the $11513$ residues,
checks that they are all nonzero, and compares the complete result
with the independently produced C++ list. The C++ calculation also
compares the original and contracted determinant formulas throughout
$12\le b\le32$, including the full exact multiplier.

The exact polynomial multiplication in the Python verifier represents
coefficients as successive base-$2^{64}$ digits of an integer. At the
largest product length here, every unreduced convolution coefficient
is smaller than $2^{64}$; hence there is no carry between coefficient
blocks. Integer multiplication is therefore precisely polynomial
convolution, with no floating-point transform or rounding. Newton
inversion is checked by multiplying the final series by
$(e^t-1)/t$ through degree $23039$ and obtaining the constant polynomial
one. Both complete replays use finite integer arithmetic. Together
with \eqref{t12:initial-bound}, they prove
Theorem~\ref{t12:complete}.

\subsection{Identity consequences and the remaining problem}

The new rank theorem yields the ordinary convergent identity family
\eqref{t10:ordinary} with the following replacements:
$M_{5,b}$ by $M_{6,b}$, $10$ by $12$, $9!$ by $11!$, and
$\psi^{(9)}$ by $\psi^{(11)}$. In the subtraction polynomial,
replace the lower outer index $b+4$ by $b+5$ and each occurrence of
$r-b-4$ by $r-b-5$. The allowed range is $b\ge7$. Thus products
of the eleventh polygamma with every higher odd-order polygamma
now reconstruct all $b$ even-zeta coordinates by explicit rational
polynomial weights and finite midpoint subtractions.

Theorems~\ref{t10:complete} and \ref{t12:complete} extend the proved
smaller orders from $2,4,6,8$ to $10,12$. The full fixed-gap-six
family and the unrestricted unequal-order conjecture remain open.
The eventual theorem reduces each further fixed smaller order to a
finite certificate, but no assertion about uncomputed orders is made
here. The leading Hankel factorization and the source contraction
are inherited; the completed order-twelve family, its simple explicit
threshold proof, and its full pair of arithmetic certificates are the
new contributions of this section.

% END sections/07_order_twelve.tex


% BEGIN sections/08_audit_and_research.tex
\section{Source reconciliation, verification, and further research}
\label{audit:section}

\subsection{Precise changes to the research status}

The targeted audit has not established an additional false theorem in
the inspected source dependencies. The proposed corrections concern
the scope and status of research questions, the distinction between
inherited and new results, and the conventions needed when combining
formulas. This distinction matters in a rapidly developing collection:
the date on an archive does not by itself determine which other results
its research list has incorporated.

\paragraph{The reconstruction questions.}
Question R1 in the latest \emph{Quartic Spectral Reconstruction}
report asks for the complete smaller-order-ten family and then the
smaller-order-twelve family \cite{tail-quartic}. Those two specified
tasks are now resolved by Theorems~\ref{t10:complete}
and~\ref{t12:complete}. In particular, the older excluded progression
for order ten is no longer an unresolved family: the inherited
eventual theorem and congruence theorem reduce the remaining task
to the single case $b=962$, which is certified here. For order twelve,
the complete finite range and the proved eventual range cover every
$b\ge7$. The unrestricted assertion for every $a<b$, requested in
R2 of the quartic report, remains open. So does the full fixed-gap-six
family $b=a+3$. The newly included cases $(p,q)=(10,16)$ and $(12,18)$
do not settle that family. Every rank statement is in the specified
formal coordinate space; no arithmetic independence of zeta values
follows from it.

\paragraph{The alternating-tail question.}
The question headed ``Bilinear adjacent-tail certificates'' in the
latest alternating-tail report \cite{tail-alternating} is answered
within exactly its two-product rational ansatz. Theorem~\ref{bi:classification}
gives necessity and sufficiency, and Theorem~\ref{bi:sum} supplies the
finite and infinite endpoint identities. The independently shifted
quadratic test case is therefore completely evaluated, including
confluence, mixed argument derivatives, and anchored rectangular
primitives. The divisor obstruction is an impossibility result for
this certificate class. It is not an impossibility theorem for all
summation identities involving an incompatible pair of denominators.
Larger support, different auxiliary terms, and independently variable
Hurwitz orders remain separate problems.

\paragraph{The Stieltjes question.}
Question 9 of the earlier \emph{Polygamma Tail Identities} report
excluded zero argument-derivative orders from its positive-order
product generator \cite{tail-polygamma}. Theorem~\ref{edge:theorem}
fills the edge with exactly one zero derivative order. Together with
the inherited positive-positive formula, the uniform dictionary
\eqref{edge:uniform-dictionary} covers every same-point quadratic
product with positive total argument-derivative order. All spectral
derivatives in that dictionary are evaluated at the single integer
$w=r+t+1$; ordinary zeta values at other positive integers occur in
the coefficients. The double-undifferentiated germ $H_{1,1}$ and the
nonzero periodic-translation formulas were already available in the
unified manuscript \cite{manuscript}. Equation~\eqref{edge:corner}
checks compatibility with the former. The translated products and
their contact corrections retain their original attribution.

The already incorporated Stieltjes source corrections, general
digamma-power identities, Gamma normal forms, and convolution towers
are not relabelled as new discoveries. The Gaussian $S_6$ and revised
$S_8$ conjectures are also unaffected. Evaluating another conductor-four
tail does not evaluate either of those separately defined periods.

\subsection{Normalization and endpoint requirements for integration}

Three conventions are essential when moving these results into the
unified manuscript. First, the finite summation identities retain their
upper endpoint until its limit is justified. In particular,
\eqref{bi:finite} and \eqref{sh:primitive-finite} record information
that can survive analytic continuation beyond the ordinary convergence
region. A continued expression is not thereby the sum of the
pointwise-specialized summands.

Second, meromorphic expressions must be combined before taking a
removable value or an order derivative. The Mellin factors involving
$s\zeta(s+1)$ are entire combinations, including at $s=0$.
At the Stieltjes edge, both resonant local terms are subtracted
jointly. Substituting a parameter equal to one into a formula proved
only for larger integer orders does not prove the extension.
Likewise, merging two separated singular points changes the local
subtraction problem. It requires a collision theorem, rather than
literal substitution into a translated finite-part formula.

Third, the finite part uses the original coordinate $x$ and the
constant term of the cutoff integral. Integration by parts therefore
has the explicit endpoint law \eqref{edge:ibp}. For example,
\[
 \FP\int_0^1\bigl\{
   \gamma_1(x)\gamma_0'(x)+\gamma_0(x)\gamma_1'(x)
   \bigr\}\,dx=-\zeta(2)-\zeta'(2),
\]
which is not zero. The Stieltjes convention
\eqref{eq:stieltjes-convention} fixes every sign in this statement.
Similarly, differentiating the weighted kernel $K(z,w)$ requires
differentiating its factor $2n+z+w$ as well as both tails. The
additional terms in \eqref{sh:mixed-formula} are part of the identity.
These are requirements for correct reuse, not assertions that the
source theorems violated them.

\subsection{What the verification actually establishes}

The order-ten certificate proves a nonzero retained determinant
$60407$ modulo $65537$, and an independent reconstruction gives
$15578$ for the corresponding original determinant modulo the same
prime. The order-twelve certificate covers all $11513$ integers
$7\le b\le11519$; its two independent implementations agree on every
residue. The infinite remaining range is handled by the written
threshold proof. These are exact arithmetic certificates. The
bilinear/kernel checks include $144$ rational cutoff identities,
$72$ rational kernel cutoff identities, and $3125$ finite divisor
strings, in addition to the symbolic cancellation and determinant
checks. Its $21$ numerical diagnostics at $50$-digit precision have
maximum absolute discrepancy below $5.4\cdot10^{-51}$.
The Mellin verifier records $2002$ exact rational or algebraic
checks and $11$ direct integer-order or scaled quadratures. The
Stieltjes verifier records $81$ exact endpoint identities, four
printed-formula checks, and five direct finite-part quadratures,
with one repeated at a second splitting point. Its largest recorded
absolute discrepancy is less than $1.3\cdot10^{-46}$ at $55$-digit
working precision. The quadratures are numerical diagnostics, not
interval certificates, and do not replace the holomorphy arguments
for arbitrary complex orders or all-index derivatives. The supplied
logs distinguish these two kinds of evidence. The combined
\texttt{verify\_all.py} run completed successfully.

\subsection{Twelve concrete questions beyond the completed results}

The following questions supplement the Mellin-specific programme
already given above. They ask for structural identities, finite
reductions, or complete nonvanishing arguments. No claim of global
novelty or of arithmetic independence is built into their formulation.

\begin{enumerate}[leftmargin=*,label=\textbf{R\arabic*.}]

\item \textbf{A uniform reconstruction mechanism.}
Find a determinant identity or a recurrence proving
$\det M_{a,b}\ne0$ for every $1\le a<b$, or determine an exact
exceptional set if one exists. The completed families through smaller
order twelve constrain any proposed formula. At present the highest
Bernoulli-degree determinant explains eventual nonvanishing for fixed
$a$, but does not exclude cancellation by the lower degrees uniformly
in both parameters. A mechanism preserving those lower-degree terms
would resolve R2 of the quartic report, rather than merely enlarging
the range of computations.

\item \textbf{The moving boundary in the gap-six family.}
Derive a bordered-Hankel or comparable exact formula for $M_{a,a+3}$
that retains all lower-index terms and the omitted-binomial
conventions at the boundary. The fixed-order contraction used in the
source does not justify omitting this information for arbitrarily
large $a$ at the fixed separation. A recurrence in $a$
with explicit nonzero initial data would be particularly useful.
This is a sharply defined intermediate problem between the existing
gap-four theorem and unrestricted unequal orders.

\item \textbf{Closed reconstruction coefficients at orders ten and twelve.}
For $a=5,6$, replace an unevaluated matrix inverse in the reconstruction
of a prescribed zeta coordinate by a uniform coefficient formula.
The target is a fixed-size kernel or product expression for the
polynomial weight, including the ordinary-integral subtraction
constants. Such a formula would turn the existence theorem for each
larger order into an immediately usable family of polygamma
integrals. Producing another determinant ratio is only useful if its
minors themselves admit a proved finite reduction.

\item \textbf{Minimal support beyond the adjacent bilinear class.}
Allow a finite collection of products $E_{n+i}F_{n+j}$ with rational
coefficients and require the difference of the certificate to contain
only a prescribed nonzero bilinear term. Determine the smallest
support needed, modulo identities already implied by the tail
recurrences. Does the divisor obstruction of
Proposition~\ref{bi:divisor} persist, or can an incompatible pair such
as generic separately shifted pure powers be accommodated?
A necessary-and-sufficient criterion would distinguish a genuine new
summation mechanism from a rewriting of the two-product certificate.

\item \textbf{Finite ranks for longer adjacent denominator ladders.}
For
\[
 L_d(z)=\sum_{j\ge0}\frac{(-1)^j}{(j+z)_d},\qquad
 \Re z>0,\quad d\ge2,
\]
determine the minimal separated rank of the weighted two-shift kernels
obtained from admissible certificates. Even $d$ admits a pairing of
adjacent linear factors; odd $d$ encounters the divisor obstruction
at generic independent shifts. The rank-two formula for $d=2$
suggests an exact finite determinant hierarchy. The objective is a
minimal rank theorem and explicit confluent determinants for every
admissible $d$, not just a separable expansion whose length has not
been proved optimal.

\item \textbf{Three independently colored tails.}
Classify finite certificates built from neighboring products of three
tails with independent denominator polynomials and colors. State
which quadratic and linear residual tails must survive and give the
complete finite endpoint. The previously developed powers of a single
tail do not settle this independent-data problem. A first substantive
target is a three-shift quadratic family, with a specified reduction
to depth-two sums and one-factor functions; unrestricted depth-one
closure should not be assumed.

\item \textbf{The continued alternating-square law at negative orders.}
Starting from its exact finite endpoint, construct a jointly
holomorphic subtraction formula at each nonpositive tail order and
derive every order derivative there. The already documented
continued value $1/4$ at order zero, despite vanishing pointwise
summands, is the compulsory normalization test
\cite{tail-alternating}. The desired identities should identify the
finite Euler-polynomial contributions and the spectral logarithmic
terms separately. This would give exact negative-order identities
without confusing continuation with an ordinary series.

\item \textbf{Independent nonperiodic Stieltjes shifts.}
Construct an order-generating formula for the ordinary integrals
\[
 \int_0^1\gamma_m^{(r)}(x+a)\gamma_n^{(t)}(x+b)\,dx,
 \qquad \Re a,\Re b>0,
\]
with $a$ and $b$ independent. These are not the inherited periodic
translations $\{x+a\}$. Begin with the antisymmetric first-derivative
combination, since its symmetric companion is an endpoint derivative.
Specify a finite function class or a genuinely new two-parameter
master, and prove how its diagonal and boundary limits recover the
appropriate coincident finite parts. The independent shift difference
is the missing variable in the classical coincident product kernel.

\item \textbf{Cubic Stieltjes products with independent indices.}
Determine an explicit holomorphic germ for
$\FP\int_0^1\gamma_m^{(r)}\gamma_n^{(t)}\gamma_j^{(v)}\,dx$
when $r+t+v>0$. A first nontrivial target is
$\FP\int_0^1\gamma_1(x)^2\gamma_0'(x)\,dx$.
The all-index digamma-power results concern a different restriction
of this problem. Fourier or Mellin analysis suggests that genuine
multiple-zeta order derivatives may be needed. A successful theorem
should identify that alphabet and all local counterterms before
claiming a finite reduction; the single-spectral-argument quadratic
closure does not imply its cubic analogue.

\item \textbf{Universal relations among the quadratic finite parts.}
For fixed total Stieltjes index and fixed positive total argument
derivative, classify all relations implied by the exact coefficient
dictionary. Determine whether exchange symmetry and the endpoint
identity generate all of them, or whether further relations remain.
The cancellation of the highest zeta derivative in
\eqref{edge:top} gives a concrete first test. Seek a uniform generating
basis and dimension in a declared formal coefficient alphabet. This
would explain cancellations systematically while avoiding an
unproved claim about independence of evaluated constants.

\item \textbf{Simultaneous and iterated triple collisions.}
For three periodically translated Stieltjes factors, compare the
constant term obtained by a simultaneous fixed-coordinate collision
with the two possible successive pairwise collisions. Compute any
difference as an explicit local distribution and test whether a
canonical correction makes the resulting composition associative.
The existing quadratic contact and dilation laws
\cite{manuscript} are mandatory compatibility conditions, not a proof
of this three-factor assertion. The question is meaningful only
after the two-variable cutoff path and the retained finite term are
specified; those choices are part of the proposed theorem.

\item \textbf{A finite evaluation of the logarithmic Gamma correction.}
The inherited function
$S(a,b)=\sum_{n\ge0}\arctan^2(b/(n+a))$ already has an ordinary
Gamma-integral representation and exact shift and distribution laws
\cite{tail-alternating}. Find a proved finite evaluation at a
nonzero mass, for example on a rational-shift locus with
$e^{-2\pi b}$ algebraic, in an explicitly stated collection of
polylogarithm values, order derivatives, and Gamma values. The
existing small-mass series and integral remain valid answers at
generic parameters. A finite special-value reduction would be an
additional identity and would clarify which extra functions the
higher spectral defects actually require.

\end{enumerate}

\subsection{Proposed repository integration}

The accompanying integration note proposes separate inserts for the
certificate classification, the shifted kernels and primitives, the
Mellin resonances, the Stieltjes edge, and the two reconstruction
families. It records their exact source dependencies and the status
sentences that should change. The current manuscript should retain
its established collision, convolution, and Gamma-normal-form
chapters, with cross-references to the added edge theorem. The
complete verification records should travel with the relevant
inserts, including the finite-field residue list and its independent
replay. No repository files are modified by this delivery.

% END sections/08_audit_and_research.tex


% BEGIN references.tex
\clearpage
\phantomsection
\addcontentsline{toc}{section}{References}
\begingroup\small
\begin{thebibliography}{99}
\raggedright
\setlength{\itemsep}{3pt}
\setlength{\parskip}{0pt}

\bibitem{manuscript}
ProveIt project,
\emph{Polylogarithms manuscript}, repository revision
\texttt{da2792fc9cc8ef19c6a314be6c65e1437fbdc66b},
11 October 2026.
\href{https://github.com/VladimirReshetnikov/ProveIt/tree/da2792fc9cc8ef19c6a314be6c65e1437fbdc66b/Analysis/Polylogarithms/docs/manuscript}{Pinned manuscript directory}.
The revision's 91 chapter files were indexed for the present audit.

\bibitem{tail-alternating}
ProveIt project,
\emph{Alternating Tails and Spectral Defects}, incoming research report,
11 October 2026. In particular, the two-product research question and
Theorem \texttt{rt:rank10} in \texttt{sections/07\_rank\_refinement.tex}.
\href{https://github.com/VladimirReshetnikov/ProveIt/blob/da2792fc9cc8ef19c6a314be6c65e1437fbdc66b/docs/incoming/ProveIt_Alternating_Tails_Spectral_Defects_2026-10-11.zip}{Pinned source archive}.

\bibitem{tail-quartic}
ProveIt project,
\emph{Quartic Spectral Reconstruction}, incoming research report,
11 October 2026. In particular, \texttt{sections/03\_eventual\_rank.tex},
Lemma \texttt{ert:normalization}, Theorem \texttt{ert:main}, the
order-five threshold certificate, and research question R1.
\href{https://github.com/VladimirReshetnikov/ProveIt/blob/da2792fc9cc8ef19c6a314be6c65e1437fbdc66b/docs/incoming/ProveIt_Quartic_Spectral_Reconstruction_2026-10-11.zip}{Pinned source archive}.

\bibitem{tail-polygamma}
ProveIt project,
\emph{Polygamma Tail Identities}, incoming research report,
11 October 2026. In particular, Theorem
\texttt{hprod:stieltjes-products} in
\texttt{sections/06\_hurwitz\_products.tex}, and paragraph 9 of
\texttt{sections/09\_audit\_research.tex}.
\href{https://github.com/VladimirReshetnikov/ProveIt/blob/446dca9bd65c20e409d1d2561d4ffad408103e2a/docs/incoming/ProveIt_Polygamma_Tail_Identities_2026-10-11.zip}{Archive at its containing commit}.
The archive's own baseline is different from that containing commit;
see the supplied provenance record.

\bibitem{Hoffman}
M. E. Hoffman,
\emph{Sums of products of Riemann zeta tails},
Mediterr. J. Math. \textbf{13} (2016), 2771--2781.
\url{https://arxiv.org/abs/1602.03399}.

\bibitem{EspinosaMoll}
O. Espinosa and V. H. Moll,
\emph{On some definite integrals involving the Hurwitz zeta function},
Ramanujan J. \textbf{6} (2002), 159--188.
\url{https://arxiv.org/abs/math/0012078}.

\bibitem{ShpotParis}
M. A. Shpot, M. P. Chaudhary, and R. B. Paris,
\emph{Integrals of products of Hurwitz zeta functions and the Casimir
 effect in $\phi^4$ field theories},
J. Class. Anal. \textbf{9} (2016), no. 2, 99--115.
DOI: \href{https://doi.org/10.7153/jca-09-11}{10.7153/jca-09-11}.
The classical Hurwitz product formula is recorded in Remark 1.
\url{https://files.ele-math.com/articles/jca-09-11.pdf}.

\bibitem{NielsenBeta}
C. Berg, S. Koumandos, and H. L. Pedersen,
\emph{Nielsen's beta function and some infinitely divisible distributions},
preprint, 2019.
\url{https://arxiv.org/abs/1905.04131}.

\bibitem{DLMFPolylog}
NIST Digital Library of Mathematical Functions,
\emph{\S25.12, Polylogarithms}, especially Eqs. 25.12.11--25.12.12.
\url{https://dlmf.nist.gov/25.12}.

\bibitem{DLMFZeta}
NIST Digital Library of Mathematical Functions,
\emph{\S25.11, Hurwitz Zeta Function}.
\url{https://dlmf.nist.gov/25.11}.

\bibitem{DLMFPolygamma}
NIST Digital Library of Mathematical Functions,
\emph{\S5.15, Polygamma Functions}.
\url{https://dlmf.nist.gov/5.15}.

\bibitem{mpr:DLMFLerch}
NIST Digital Library of Mathematical Functions,
\emph{\S25.14, Lerch's Transcendent}, especially the defining series and
integral representations.
\url{https://dlmf.nist.gov/25.14}.

\bibitem{mpr:DLMFbeta}
NIST Digital Library of Mathematical Functions,
\emph{\S5.12, Beta Function}.
\url{https://dlmf.nist.gov/5.12}.

\bibitem{mpr:Freitas}
P.~Freitas,
\emph{Integrals of polylogarithmic functions, recurrence relations, and
associated Euler sums},
Math. Comp. \textbf{74} (2005), 1425--1440.
Preprint: \url{https://arxiv.org/abs/math/0406401}.

\bibitem{mpr:Sofo}
A.~Sofo,
\emph{Integrals of polylogarithmic functions with negative argument},
Acta Univ. Sapientiae Math. \textbf{10} (2018), no.~2, 347--367.
\url{https://doi.org/10.2478/ausm-2018-0027}.
Publisher PDF archived by the author's institution:
\url{https://vuir.vu.edu.au/45960/1/10.2478_ausm-2018-0027.pdf}.

\end{thebibliography}
\endgroup

% END references.tex

\end{document}
